\documentclass[11pt,reqno]{amsart}

\usepackage[T1]{fontenc}
\usepackage[utf8]{inputenc}
\usepackage{lmodern}
\usepackage{microtype}
\usepackage{amsmath,amssymb,amsthm,mathtools,mathrsfs}
\usepackage{enumitem}
\usepackage{booktabs}
\usepackage{longtable}
\usepackage{array}
\usepackage{xcolor}
\usepackage{hyperref}
\usepackage[nameinlink,capitalize,noabbrev]{cleveref}
\usepackage{aliascnt}
\usepackage{tikz}
\usetikzlibrary{arrows.meta,positioning}

\hypersetup{
  colorlinks=true,
  linkcolor=blue!45!black,
  citecolor=green!35!black,
  urlcolor=blue!55!black,
  pdftitle={Separator algebras and transfer theory for the zero-forcing polynomial},
  pdfauthor={Guillaume Sidoine}
}

\newtheorem{theorem}{Theorem}[section]
\newaliascnt{lemma}{theorem}
\newtheorem{lemma}[lemma]{Lemma}
\aliascntresetthe{lemma}
\newaliascnt{proposition}{theorem}
\newtheorem{proposition}[proposition]{Proposition}
\aliascntresetthe{proposition}
\newaliascnt{corollary}{theorem}
\newtheorem{corollary}[corollary]{Corollary}
\aliascntresetthe{corollary}
\newaliascnt{claim}{theorem}
\newtheorem{claim}[claim]{Claim}
\aliascntresetthe{claim}
\theoremstyle{definition}
\newaliascnt{definition}{theorem}
\newtheorem{definition}[definition]{Definition}
\aliascntresetthe{definition}
\newaliascnt{example}{theorem}
\newtheorem{example}[example]{Example}
\aliascntresetthe{example}
\newaliascnt{remark}{theorem}
\newtheorem{remark}[remark]{Remark}
\aliascntresetthe{remark}

\crefname{theorem}{Theorem}{Theorems}
\Crefname{theorem}{Theorem}{Theorems}
\crefname{lemma}{Lemma}{Lemmas}
\Crefname{lemma}{Lemma}{Lemmas}
\crefname{proposition}{Proposition}{Propositions}
\Crefname{proposition}{Proposition}{Propositions}
\crefname{corollary}{Corollary}{Corollaries}
\Crefname{corollary}{Corollary}{Corollaries}
\crefname{definition}{Definition}{Definitions}
\Crefname{definition}{Definition}{Definitions}
\crefname{example}{Example}{Examples}
\Crefname{example}{Example}{Examples}
\crefname{remark}{Remark}{Remarks}
\Crefname{remark}{Remark}{Remarks}

\newcommand{\Zpoly}{\mathcal Z}
\newcommand{\cl}{\operatorname{cl}}
\newcommand{\Pstate}{\mathsf P}
\newcommand{\Astate}{\mathsf A}
\newcommand{\Dalg}{\mathbb D}
\newcommand{\N}{\mathbb N}
\newcommand{\eps}{\varepsilon}
\newcommand{\one}{\mathbf 1}

\title[Separator algebras for the zero-forcing polynomial]
{Separator algebras and transfer theory\\
for the zero-forcing polynomial}

\author{Guillaume Sidoine}
\address{Independent researcher, Mauritius}
\thanks{ORCID: 0009-0005-4525-1182.}

\subjclass[2020]{05C31, 05C50, 05C76, 05C85, 68R10}
\keywords{zero forcing, graph polynomial, separator, connection matrix, graph algebra, fort, transfer algebra, block decomposition}

\begin{document}

\begin{abstract}
For a finite graph $G$, let
\[
  \Zpoly(G;x)=\sum_k z(G;k)x^k,
\]
where $z(G;k)$ is the number of zero-forcing sets of size $k$.  Boyer et al.
asked for splitting formulas for this polynomial over cut vertices and
separating sets.  We give such a formula for every finite labeled separator.

The construction uses forts.  For a separator of order $q$, a local obstruction
is encoded by a join-closed family of words in
$\{0,1,\mathsf I\}^q\setminus\{0^q\}$.  We give the exact rule for gluing two
local families and express the boundary choice as a transversal problem on the
resulting fort traces.  The corresponding connection kernels admit a M\"obius
compression indexed by chains.  Every chain state is realized by a finite
simple graph, so the finite chain kernel has the same rank as the infinite
zero-forcing connection matrix.

For an independent boundary this rank is the number $c_q$ of chains in
$C_3^q\setminus\{0^q\}$.  In particular
$c_1=4$, $c_2=52$, $c_3=1636$, and
$\log c_q=2q\log q+O(q)$.  With arbitrary boundary edges the first full ranks
are $r_0=1$, $r_1=4$, and $r_2=60$, the last value certified by an exact finite
calculation.  At one terminal the theory reduces to
$\mathbb Q(x)[s,a]/(s^2-s,a^2)$; the associated transfers give explicit
quadratic coefficient-operation algorithms for trees and block-cactus graphs.
\end{abstract}

\maketitle

\section{Introduction}

Zero forcing starts from a set of blue vertices.  A blue vertex may force its
unique white neighbour, and an initial set is zero forcing when repeated legal
forces colour the whole graph.  The parameter arose from minimum-rank and
maximum-nullity questions for symmetric matrices described by a graph
\cite{AIM2008}.

Boyer, Brimkov, English, Ferrero, Keller, Kirsch, Phillips, and Reinhart
\cite{BoyerEtAl2019} introduced
\[
  \Zpoly(G;x)=\sum_{k=0}^{|V(G)|}z(G;k)x^k,
\]
where $z(G;k)$ counts zero-forcing sets of size $k$.  They asked, among other
things, for splitting formulas based on cut vertices or separating sets.  We
answer that question for separators of arbitrary finite order.

We use connection matrices to state the answer.  Two $q$-boundaried graphs are
glued by identifying vertices with the same boundary label, and the entry of
the connection matrix is the zero-forcing polynomial of the glued graph.  The
rank records how much linear information must pass through the boundary.  This
language is standard in graph algebras \cite{Lovasz2006,KotekMakowsky2014}; in
particular, general finite-rank and bilinear-reduction theorems are already
known for suitable definable graph polynomials.  Our point is more specific:
we give the zero-forcing interface explicitly, reduce it to a finite kernel,
and determine that kernel exactly enough to recover the connection rank.

The static obstruction is a fort.  A nonempty set $F$ is a fort if no vertex
outside $F$ has exactly one neighbour in $F$, and a set is zero forcing if and
only if it meets every fort \cite{AIM2008}.  The fort hypergraph and its
transversals have been studied in their own right
\cite{CameronHogbenEtAl2026}.  Recent work also counts inclusion-minimal forts,
including the linear-time tree algorithm of Abiad and Ghasemi Nezhad
\cite{BeckerEtAl2025,CameronLi2025,DatKenter2026,AbiadGhasemiNezhad2026}.
That is a different counting problem.  Here the object being counted is every
subset that meets every fort, with its cardinality retained.

Let $B$ be a separator of order $q$.  We impose the fort condition away from
$B$ and call the resulting object a prefort.  At a boundary vertex only three
pieces of information matter: the vertex lies in the prefort, it lies outside
and has no interior prefort neighbour, or it lies outside and has exactly one.
Counts at least two are already safe.  The resulting atoms lie in
\[
  P_q=\{0,1,\mathsf I\}^q\setminus\{0^q\}.
\]
For a fixed interior initial set the atom family is join closed in the product
order $0<1<\mathsf I$.  If two atoms have the same trace $T$ and the glued
boundary has edge set $E$, they form a fort exactly when
\[
  d_E(i,T)+\alpha_i+\beta_i\ne1\qquad(i\notin T).
\]
The selected boundary vertices must then meet every trace produced by a
compatible pair.  This is the splitting formula.

The local state space is large, but its linear content is smaller.  A M\"obius
cancellation on finite join-semilattices leaves only chains.  We prove that
every chain state is realized by a finite simple graph and use one-terminal
selection projectors to isolate its prescribed interior set.  Hence no mode in
the compressed kernel is formal only.  If $M_E(x)$ denotes the chain kernel
for boundary edge set $E$, then
\[
  r_q=\sum_{R\subseteq\binom{[q]}2}
  \operatorname{rank}_{\mathbb Q(x)}
  \left(\sum_{E\supseteq R}(-1)^{|E\setminus R|}M_E(x)\right).
\]
For an independent boundary the rank is the number $c_q$ of chains in
$C_3^q\setminus\{0^q\}$.  We obtain an explicit formula for $c_q$ and
\[
 c_0=1,\qquad c_1=4,\qquad c_2=52,\qquad c_3=1636,
 \qquad \log c_q=2q\log q+O(q).
\]
With arbitrary boundary edges, the first nontrivial full value beyond one
terminal is $r_2=60$; the finite rank calculation is supplied with an
independent certificate and verifier.

At $q=1$ the general construction collapses to four responses: a rooted piece
may or may not supply the root, and may or may not need the root to spend its
single force inside the piece.  The resulting algebra is
\[
  \mathbb Q(x)[s,a]/(s^2-s,a^2).
\]
This small case is useful computationally.  It gives an exact edge transfer,
a quadratic tree algorithm, block transfers, a linear clique transfer, and a
finite cycle recognizer.  For block-cactus graphs the full polynomial is
computed in $O(n^2)$ coefficient operations.  General tractability of that
class also follows from bounded clique-width and MSO machinery; the value of
the transfer here is the explicit state system and its concrete arithmetic
bound.

The first part of the paper develops the rooted and separator algebras.  The
second gives the explicit tree, block, clique, and cycle transfers.  Finite
computations used for $r_2$ and for the reduced cycle automaton are separated
from the symbolic proofs and supplied with independent verifiers.

\section{Preliminaries and restricted closure}\label{sec:preliminaries}

All graphs are finite and simple.  Let $G=(V,E)$ and let $S\subseteq V$.  Vertices of $S$ are initially blue and the remaining vertices are white.  The \emph{zero-forcing colour-change rule} is
\[
  u\to w
  \quad\text{if $u$ is blue and $w$ is the unique white neighbour of $u$.}
\]
A set $S$ is a \emph{zero-forcing set} if a sequence of legal forces colours every vertex.  We write $z(G;k)$ for the number of zero-forcing sets of cardinality $k$, and
\[
  \Zpoly(G;x)=\sum_{k=0}^{|V|}z(G;k)x^k.
\]
For a non-empty graph the constant coefficient is zero, but allowing $k=0$ is convenient for products and for the empty graph.

We shall repeatedly use the elementary observation that a vertex performs at most one force.

\begin{lemma}\label{lem:one-force}
Every vertex is the forcing vertex in at most one step of a zero-forcing process.
\end{lemma}

\begin{proof}
If $u$ forces $w$, then immediately after the force every neighbour of $u$ is blue.  Since blue vertices never become white, $u$ can never again have a white neighbour.
\end{proof}

Let $(H,r)$ be a rooted graph.  We require a closure operation in which the root is prohibited from acting.

\begin{definition}
Let $Q\subseteq V(H)$.  For $X\subseteq V(H)$, the \emph{$Q$-restricted closure} of $X$, denoted
\[
  \cl_H^{\neg Q}(X),
\]
is the terminal blue set obtained by exhausting legal forces while no vertex of $Q$ is permitted to act as a forcing vertex.  Vertices of $Q$ may nevertheless be forced.  For a singleton $Q=\{r\}$ we write $\cl_H^{\neg r}(X)$.
\end{definition}

The next lemma justifies the notation.

\begin{lemma}[Confluence and monotonicity]\label{lem:closure}
The set $\cl_H^{\neg Q}(X)$ is independent of the order in which legal $Q$-restricted forces are applied.  Moreover,
\[
  X\subseteq Y
  \quad\Longrightarrow\quad
  \cl_H^{\neg Q}(X)\subseteq \cl_H^{\neg Q}(Y).
\]
\end{lemma}

\begin{proof}
For $X\subseteq V(H)$ define
\[
  F_Q(X)=X\cup
  \left\{
    w:\text{there exists }u\in X\setminus Q
    \text{ with }N_H(u)\setminus X=\{w\}
  \right\}.
\]
The map $F_Q$ is monotone.  Indeed, let $X\subseteq Y$ and suppose that $w\in F_Q(X)\setminus Y$.  Then some $u\in X\setminus Q$ has $w$ as its unique neighbour outside $X$.  Since $w\notin Y$, enlarging $X$ to $Y$ does not remove $w$ and cannot introduce a new neighbour outside the blue set.  Hence $w$ remains the unique neighbour of $u$ outside $Y$, so $w\in F_Q(Y)$.

Since $V(H)$ is finite, iteration of $F_Q$ from $X$ reaches the least fixed point of $F_Q$ containing $X$.  Every legal sequential $Q$-restricted forcing process remains inside this fixed point.  Conversely, the vertices introduced at each synchronous iterate may be added in any order by legal $Q$-restricted forces.  The terminal set is therefore unique, and monotonicity follows by iterating the monotone map $F_Q$.
\end{proof}

For $S\subseteq V(H)\setminus\{r\}$ put
\[
  C^-(S)=\cl_H^{\neg r}(S),
  \qquad
  C^+(S)=\cl_H^{\neg r}(S\cup\{r\}).
\]
The superscript records whether the root is initially white or blue.

\begin{lemma}[Monotone simulation of a forcing sequence]\label{lem:sequence-simulation}
Let $B\subseteq B'$ be blue sets in a graph, and let
\[
  u_1\to w_1,\ldots,u_t\to w_t
\]
be a legal forcing sequence from $B$, with all forcing vertices outside a fixed
forbidden set $Q$.  Starting from $B'$, scan the same list in order and perform
$u_i\to w_i$ whenever $w_i$ is still white, omitting the step otherwise.  Every
performed step is legal, no forcing vertex lies in $Q$, and the final blue set
contains the final blue set obtained from $B$.
\end{lemma}

\begin{proof}
Induct on the list.  At a step $u\to w$, if $w$ is already blue there is
nothing to prove.  Otherwise the current larger blue set contains the current
smaller one.  The vertex $u$ has no white neighbour in the smaller process
other than $w$; enlarging the blue set cannot create a new white neighbour.
The vertex $w$ remains the unique white neighbour of $u$, so the step is legal.
\end{proof}

\begin{lemma}[Source-closure identity]\label{lem:source-closure}
If $r\in C^-(S)$, then
\[
  C^-(S)=C^+(S).
\]
\end{lemma}

\begin{proof}
Monotonicity gives $C^-(S)\subseteq C^+(S)$.  The set $C^-(S)$ is an
$r$-restricted fixed point and, because it contains $r$, it contains
$S\cup\{r\}$.  Since $C^+(S)$ is the least such fixed point, the reverse
inclusion follows.
\end{proof}

\section{The four rooted response states}\label{sec:states}

Fix a rooted graph $(H,r)$ and a set $S\subseteq V(H)\setminus\{r\}$.

\begin{definition}[Source indicator]
The source indicator of $S$ is
\[
  \sigma_r(S)=
  \begin{cases}
    1,&r\in C^-(S),\\
    0,&r\notin C^-(S).
  \end{cases}
\]
Accordingly, $\sigma_r(S)=1$ precisely when the non-root vertices can force an initially white root without using the root as a forcing vertex.
\end{definition}

We next classify the successful behaviour when the root is initially blue.

\begin{definition}[Passive, active, and blocked states]\label{def:states}
The set $S$ is called
\begin{enumerate}[label=\textnormal{(\roman*)}]
\item \emph{passive} if $C^+(S)=V(H)$;
\item \emph{active} if $C^+(S)\ne V(H)$ and there is a vertex $u$ such that
  \[
    N_H(r)\setminus C^+(S)=\{u\}
  \]
  and
  \[
    \cl_H^{\neg r}\bigl(C^+(S)\cup\{u\}\bigr)=V(H);
  \]
\item \emph{blocked} otherwise.
\end{enumerate}
\end{definition}

A passive set completes the graph without a force by the root.  An active set completes the graph after the root performs its unique available force.  A blocked set cannot be completed even if the root is initially blue.

For $\eta\in\{0,1\}$ define
\begin{align*}
  P_\eta(H,r;x)
  &=\sum_{\substack{S\subseteq V(H)\setminus\{r\}\\
                     S\text{ passive},\ \sigma_r(S)=\eta}}
      x^{|S|},\\
  A_\eta(H,r;x)
  &=\sum_{\substack{S\subseteq V(H)\setminus\{r\}\\
                     S\text{ active},\ \sigma_r(S)=\eta}}
      x^{|S|}.
\end{align*}
We suppress $(H,r;x)$ when no ambiguity can arise, and write
\[
  P=P_0+P_1,
  \qquad
  A=A_0+A_1.
\]
Blocked sets will not enter the transfer algebra because they cannot occur in a successful global configuration.

\section{Root normal form}\label{sec:normal-form}

The following normalization statement makes precise the reordering used
throughout the paper.

\begin{lemma}[Root normalization]\label{lem:root-normalization}
Let $(H,r)$ be rooted, let $S\subseteq V(H)\setminus\{r\}$, and suppose that
$r$ is initially blue.  Every successful forcing sequence can be replaced by
one of the form
\[
  \text{$r$-restricted saturation}
  \ ;\ 
  \text{at most one force by $r$}
  \ ;\ 
  \text{$r$-restricted saturation}.
\]
More precisely, either $C^+(S)=V(H)$, or there is a unique vertex
$u\in N_H(r)\setminus C^+(S)$ such that $r\to u$ is legal from $C^+(S)$ and
\[
  \cl_H^{\neg r}\bigl(C^+(S)\cup\{u\}\bigr)=V(H).
\]
\end{lemma}

\begin{proof}
Fix a successful sequence.  If the root is never used, the entire sequence is
$r$-restricted, and confluence gives $C^+(S)=V(H)$.

Otherwise let $r\to u$ be the (necessarily unique) force by $r$.  Replace the
prefix preceding this step by any restricted sequence which saturates
$S\cup\{r\}$ to $C^+(S)$.  The blue set has only increased.  If $u$ is still
white, every other neighbour of $r$ was blue before the original force and is
therefore blue in $C^+(S)$, so $r\to u$ remains legal.  If $u$ has already
become blue, omit the force by $r$.  In either case, apply
\cref{lem:sequence-simulation} to the remaining suffix.  Since $r$ cannot
force a second time by \cref{lem:one-force}, every retained force in the
suffix is restricted.  The resulting sequence has the asserted form and
still colours all vertices.
\end{proof}

\begin{theorem}[Root normal form]\label{thm:normal-form}
Let $(H,r)$ be a rooted graph and let $S\subseteq V(H)\setminus\{r\}$.
\begin{enumerate}[label=\textnormal{(\alph*)}]
\item The set $S\cup\{r\}$ is zero forcing in $H$ if and only if $S$ is passive or active.
\item The set $S$ is zero forcing in $H$ if and only if $S$ is passive or active and $\sigma_r(S)=1$.
\end{enumerate}
Consequently,
\begin{equation}\label{eq:normal-form}
  \boxed{
  \Zpoly(H;x)=(1+x)(P+A)-(P_0+A_0).
  }
\end{equation}
\end{theorem}

\begin{proof}
Part~(a) is exactly \cref{lem:root-normalization} together with
\cref{def:states}.  For part~(b), a successful process with $r$ initially
white must first colour $r$ by a force whose forcing vertex is not $r$; hence
$r\in C^-(S)$.  Conversely, if $\sigma_r(S)=1$, execute a restricted sequence
which colours $r$.  By \cref{lem:source-closure}, the resulting restricted
closure is $C^+(S)$, after which the passive or active completion from
part~(a) applies.

Sets containing $r$ contribute $x(P+A)$, whereas sets excluding $r$ contribute
$P_1+A_1=(P+A)-(P_0+A_0)$.  This proves \eqref{eq:normal-form}.
\end{proof}

\begin{remark}
The four states record precisely the two pieces of information which can cross
a one-vertex boundary: whether the rooted component consumes the root's unique
force, and whether it can supply the root when the root is initially white.
\end{remark}

\section{Rooted one-point unions and dual numbers}\label{sec:wedge}

Let $(H,r)$ and $(K,s)$ be rooted graphs with otherwise disjoint vertex sets.  Their \emph{rooted one-point union}, or \emph{wedge}, is the rooted graph
\[
  (H,r)\vee(K,s)
\]
obtained by identifying $r$ and $s$; the identified vertex is the root.

\begin{proposition}[Set-level wedge composition]\label{prop:wedge-set}
Let $S_H\subseteq V(H)\setminus\{r\}$ and
$S_K\subseteq V(K)\setminus\{s\}$, and identify $r=s$ to a common root $v$.
Then the following hold.
\begin{enumerate}[label=\textnormal{(\roman*)}]
\item The wedge is passive exactly when both component sets are passive.
\item It is active exactly when one component set is active and the other is passive.
\item Every other pair of component states is blocked.
\item Its source indicator is the Boolean disjunction of the two component source indicators.
\end{enumerate}
\end{proposition}

\begin{proof}
With $v$ initially blue and forbidden from forcing, the non-root vertices in
$H$ and $K$ have no neighbours across the separator.  Thus their restricted
closures are computed independently.  A passive component completes without
using $v$; an active component leaves one internal white neighbour which only
$v$ can initially force; and a blocked component cannot be completed even
when $v$ is supplied and allowed its one force.  It follows that a blocked
component cannot be rescued by activity on the other side, that two active
components cannot both be completed, and that exactly one active component
can be completed when every other component is passive.  This proves
(i)--(iii), using \cref{lem:one-force}.

For (iv), if neither side can force the initially white root, then no first
force colouring $v$ exists in the wedge.  If one side can source $v$, execute
its restricted source sequence first.  By \cref{lem:source-closure}, that side
then has exactly the same restricted closure as if $v$ had been supplied at
time zero.  The other side may therefore proceed from its plus closure.  In
particular, a passive source branch may supply $v$ before an active branch
consumes the force of $v$, and an active source branch may first supply $v$
and then consume that force internally.  No additional timing state is
needed.
\end{proof}

\begin{theorem}[Rooted product]\label{thm:wedge}
Let
\[
  (P^H,P_0^H,A^H,A_0^H)
  \quad\text{and}\quad
  (P^K,P_0^K,A^K,A_0^K)
\]
be the response data of two rooted graphs.  The response data of their rooted
wedge satisfy
\begin{align}
  P&=P^HP^K, & P_0&=P_0^HP_0^K,\label{eq:wedge-passive}\\
  A&=A^HP^K+P^HA^K,
  & A_0&=A_0^HP_0^K+P_0^HA_0^K.\label{eq:wedge-active}
\end{align}
\end{theorem}

\begin{proof}
Apply \cref{prop:wedge-set} to every ordered pair of component subsets and
sum their generating weights.  The source-zero terms are precisely those in
which both component source indicators are zero.
\end{proof}

The product becomes particularly transparent over the algebra of dual numbers.  Put
\[
  \Dalg=\mathbb Z[x,\eps]/(\eps^2)
\]
and define
\begin{align*}
  \mathscr R(H,r)&=P(H,r)+\eps A(H,r),\\
  \mathscr R_0(H,r)&=P_0(H,r)+\eps A_0(H,r).
\end{align*}

\begin{corollary}[Dual-number multiplicativity]\label{cor:dual}
For rooted graphs $(H,r)$ and $(K,s)$,
\begin{align*}
  \mathscr R\bigl((H,r)\vee(K,s)\bigr)
  &=\mathscr R(H,r)\mathscr R(K,s),\\
  \mathscr R_0\bigl((H,r)\vee(K,s)\bigr)
  &=\mathscr R_0(H,r)\mathscr R_0(K,s).
\end{align*}
\end{corollary}

\begin{proof}
This is exactly \cref{thm:wedge}, since the coefficient of $\eps^2$ vanishes.
\end{proof}

The four coordinates are not merely sufficient; in the natural connection-matrix
sense they are necessary.

\begin{theorem}[Contextual completeness and rank four]\label{thm:connection-rank}
Work over $K=\mathbb Q(x)$ and write
\[
  r(H)=(P_0(H),P_1(H),A_0(H),A_1(H))^T
\]
for a finite rooted graph $H$.  Let $C$ be the connection matrix indexed by
finite rooted graphs, with
\[
  C_{H,L}=\Zpoly(H\vee L;x).
\]
Then $\operatorname{rank}_K C=4$.  Moreover, two rooted graphs have the same
response profile if and only if
\[
  \Zpoly(H\vee L;x)=\Zpoly(H'\vee L;x)
\]
for every rooted graph $L$.  It is enough to test the four fixed contexts
$K_1$, $K_2$ rooted at an endpoint, and $P_3$ rooted respectively at an
endpoint and at its middle vertex.
\end{theorem}

\begin{proof}
By \cref{prop:wedge-set,thm:normal-form}, compatibility of two successful
states is bilinear.  In the coordinate order $(P_0,P_1,A_0,A_1)$,
\[
  C_{H,L}=r(H)^T B(x)r(L),
\]
where
\[
B(x)=
\begin{pmatrix}
 x&x+1&x&x+1\\
 x+1&x+1&x+1&x+1\\
 x&x+1&0&0\\
 x+1&x+1&0&0
\end{pmatrix}.
\]
Indeed, an incompatible pair of active states contributes zero; a compatible
pair contributes $x$ when the common root must be initially selected, and an
additional $1$ exactly when at least one side can source an initially white
root.  Hence $\operatorname{rank}C\le4$.

For the four rooted contexts listed in the statement, direct use of the
profiles in \cref{subsec:sanity} gives the row matrix
\[
V(x)=
\begin{pmatrix}
 1&0&0&0\\
 0&x&1&0\\
 x&x+x^2&1&0\\
 0&x^2&0&2x
\end{pmatrix}.
\]
A direct calculation gives
\[
  \det V=-2x^3,
  \qquad
  \det B=(x+1)^2.
\]
The corresponding $4\times4$ connection minor is $VBV^T$ and therefore has
\[
  \det(VBV^T)=4x^6(x+1)^2\ne0.
\]
Thus $\operatorname{rank}C=4$.  If the difference of two response vectors
pairs to zero with the four displayed contexts, invertibility of $B$ and $V$
forces that difference to vanish.  The converse is immediate from the
factorization.
\end{proof}

\begin{remark}[Full response algebra]\label{rem:full-algebra}
The complete set-level multiplication can be encoded over $K$ by
\[
  P_0\leftrightarrow1,\qquad P_1\leftrightarrow s,
  \qquad A_0\leftrightarrow a,\qquad A_1\leftrightarrow sa,
\]
with
\[
  K[s,a]/(s^2-s,a^2)
  \cong K[a]/(a^2)\times K[a]/(a^2).
\]
The idempotent $s$ records Boolean source disjunction and the nilpotent $a$
records consumption of the root's unique force.  Evaluating the source
coordinate at $s=0$ and $s=1$ recovers respectively the source-zero and total
dual-number profiles $\mathscr R_0$ and $\mathscr R$.  Thus the dual-number
notation is a pair of projections of the full four-state algebra.
\end{remark}

\begin{corollary}[Scalar specializations]\label{cor:connection-specialization}
For the numerical invariant $\Zpoly(G;t)$ over a characteristic-zero field,
the rooted connection matrix has rank four when $t\notin\{0,-1\}$, rank two
when $t=-1$, and rank zero when $t=0$.
\end{corollary}

\begin{proof}
For $t\ne0$, the displayed matrix $V(t)$ is invertible, so the rank equals that
of $B(t)$.  This is four except at $t=-1$, where the nonzero part of $B(-1)$
has rank two.  At $t=0$, every rooted wedge is nonempty and its zero-forcing
polynomial has zero constant coefficient.
\end{proof}

For a wedge of $d$ rooted graphs the response is therefore
\begin{align}
  P&=\prod_{i=1}^dP_i,\label{eq:dpassive}\\
  A&=\sum_{j=1}^d A_j\prod_{i\ne j}P_i,\label{eq:dactive}
\end{align}
with the analogous formulas for $P_0$ and $A_0$.

\subsection{An exact cut-vertex formula}

Let $G$ be connected and let $v$ be a cut vertex.  Let $W_1,\dots,W_d$ be the components of $G-v$, and put
\[
  B_i=G[W_i\cup\{v\}],
\]
rooted at $v$.  For each branch let $P_i,P_{i,0},A_i,A_{i,0}$ denote its response polynomials.

\begin{theorem}[Articulation formula]\label{thm:articulation}
Define
\begin{align*}
  R_v(G;x)
  &=\prod_{i=1}^dP_i
    +\sum_{j=1}^dA_j\prod_{i\ne j}P_i,\\
  R_{v,0}(G;x)
  &=\prod_{i=1}^dP_{i,0}
    +\sum_{j=1}^dA_{j,0}\prod_{i\ne j}P_{i,0}.
\end{align*}
Then
\begin{equation}\label{eq:articulation}
  \boxed{
  \Zpoly(G;x)=(1+x)R_v(G;x)-R_{v,0}(G;x).
  }
\end{equation}
The contribution of zero-forcing sets containing $v$ is $xR_v(G;x)$, and the contribution of those excluding $v$ is $R_v(G;x)-R_{v,0}(G;x)$.
\end{theorem}

\begin{proof}
The rooted graph $(G,v)$ is the wedge of the branches $(B_i,v)$.  Apply \cref{thm:wedge,thm:normal-form}.
\end{proof}

\begin{remark}
Row's cut-vertex method \cite{Row2012} determines the least cardinality of a zero-forcing set by tracking zero spread.  Formula \eqref{eq:articulation} is coefficientwise: it reconstructs the number of zero-forcing sets of every cardinality.
\end{remark}

\section{Arbitrary separators and atomic preforts}\label{sec:separator-preforts}

We now pass from a one-vertex boundary to an arbitrary finite separator.  The
one-terminal response states remain useful later as selection projectors, but
the natural general interface is most transparent in the fort language.
Throughout this section $B=[q]=\{1,\ldots,q\}$ is an ordered boundary and the
corresponding vertices of a boundaried graph are denoted $b_1,\ldots,b_q$.

\begin{definition}[Boundaried graph and gluing]\label{def:q-boundaried}
A \emph{$q$-boundaried graph} is a finite simple graph $H$ together with
pairwise distinct labeled vertices $b_1,\ldots,b_q$.  If $H$ and $K$ are
$q$-boundaried graphs, $H\oplus_q K$ is obtained from their disjoint union by
identifying equally labeled boundary vertices and suppressing duplicate edges.
Let $E_B(H)\subseteq\binom{B}{2}$ denote the graph induced by the boundary.
\end{definition}

The $q$th connection matrix of the zero-forcing polynomial is the infinite
matrix
\[
  C_q[H,K]=\Zpoly(H\oplus_q K;x)
\]
over all isomorphism classes of $q$-boundaried graphs.  We write
\[
  r_q=\operatorname{rank}_{\mathbb Q(x)}C_q.
\]
Unless a specialization is stated explicitly, every connection rank in this
paper is taken over the rational-function field $\mathbb Q(x)$.
For $q=0$, multiplicativity on disjoint unions gives $r_0=1$.  Theorem
\ref{thm:connection-rank} gives $r_1=4$.

\subsection{Preforts and atomic words}

Recall that a fort is a nonempty set $F$ such that no vertex outside $F$ has
exactly one neighbour in $F$.  Equivalently, a set is zero forcing if and only
if it meets every fort \cite{AIM2008}.  We omit the boundary conditions locally.

\begin{definition}[Prefort]\label{def:prefort}
Let $H$ be $q$-boundaried and let
$S\subseteq V(H)\setminus B$ be an interior initial set.  A nonempty set
$F\subseteq V(H)\setminus S$ is a \emph{$B$-prefort} if
\[
 |N_H(v)\cap F|\ne1
 \qquad\text{for every }v\in V(H)\setminus(F\cup B).
\]
The fort condition is imposed at every nonboundary vertex and omitted at
$B$.
\end{definition}

Write $\mathsf I$ for the symbol recording boundary membership.  For a prefort
$F$ and $i\in B$ put
\[
 c_i(F)=|N_H(b_i)\cap(F\setminus B)|.
\]
Its set of \emph{atomic refinements} consists of all words
$\alpha\in\{0,1,\mathsf I\}^q$ whose $i$th coordinate is
\[
 \alpha_i=\begin{cases}
 \mathsf I,&b_i\in F,\\
 0,&b_i\notin F\text{ and }c_i(F)=0,\\
 1,&b_i\notin F\text{ and }c_i(F)=1,\\
 0\text{ or }1,&b_i\notin F\text{ and }c_i(F)\ge2.
 \end{cases}
\]
The last line is the crucial saturation rule: once two interior fort
neighbours are already present, adding an external contribution can never
make the final count equal to one, so either binary refinement is safe for an
existential compatibility test.

Let
\[
  \mathcal A_H(S)
\]
be the union of the atomic refinements of all $B$-preforts disjoint from $S$.
The all-zero atom has a special status.

\begin{lemma}[Fatal atom]\label{lem:fatal-atom}
If $0^q\in\mathcal A_H(S)$, then $S$ cannot occur as the interior part of a
zero-forcing set in $H\oplus_q K$ for any external $q$-boundaried graph $K$ and
any choice of initially blue boundary vertices.
\end{lemma}

\begin{proof}
The atom $0^q$ is supplied by a nonempty prefort $F$ containing no boundary
vertex and for which every boundary vertex has either zero or at least two
neighbours in $F$.  Hence the omitted fort conditions at the boundary are
already satisfied.  It follows that $F$ is a fort of $H$, and it remains a fort after any
external graph is glued along $B$, because no new vertex of $F$ is introduced
and external vertices have no neighbours in $F$ except through boundary
vertices, none of which lies in $F$.  Since $F\cap S=\varnothing$, no extension
of $S$ is zero forcing.
\end{proof}

We therefore work with the nonfatal atom poset
\[
  P_q=\{0,1,\mathsf I\}^q\setminus\{0^q\},
\]
ordered coordinatewise by
\[
  0<1<\mathsf I.
\]
The join $\alpha\vee\beta$ is the coordinatewise maximum.

\begin{lemma}[Join closure]\label{lem:atom-join-closed}
If $S$ is nonfatal, then $\mathcal A_H(S)\subseteq P_q$ is join closed:
\[
 \alpha,\beta\in\mathcal A_H(S)
 \quad\Longrightarrow\quad
 \alpha\vee\beta\in\mathcal A_H(S).
\]
\end{lemma}

\begin{proof}
Choose preforts $F,G$ and atomic refinements producing $\alpha,\beta$.
The union $F\cup G$ is again a prefort.  Indeed, if a nonboundary vertex lies
outside $F\cup G$, then its number of neighbours in each of $F$ and $G$ is
zero or at least two; the same is therefore true for the union.

At a boundary coordinate, membership in either prefort gives the symbol
$\mathsf I$ in the union.  If neither contains the boundary vertex and at
least one chosen atom has value $1$, then the union has at least one interior
prefort neighbour; if the actual count is one the refinement is forced to be
$1$, and if it is at least two the refinement $1$ is available.  If both
chosen values are zero, the union either has count zero or at least two, in
both cases allowing the refinement zero.  The union $F\cup G$ supplies
$\alpha\vee\beta$.
\end{proof}

Let $\mathfrak J_q$ denote the finite set of join-closed subsets of $P_q$,
including the empty set.  Every nonfatal decorated piece $(H,S)$ therefore has
a state
\[
  \Theta(H,S)=\bigl(E_B(H),\mathcal A_H(S)\bigr)
  \in 2^{\binom B2}\times\mathfrak J_q.
\]

\subsection{The exact gluing criterion}

For an atom $\alpha$ define its trace
\[
 T(\alpha)=\{i\in B:\alpha_i=\mathsf I\}.
\]
Let $E\subseteq\binom B2$ be a boundary-edge graph and put
\[
 d_E(i,T)=|N_E(i)\cap T|.
\]

\begin{definition}[Atomic compatibility]\label{def:atomic-compatible}
Atoms $\alpha,\beta\in P_q$ are \emph{$E$-compatible} if they have the same
trace $T$ and
\begin{equation}\label{eq:atomic-compatibility}
 d_E(i,T)+\alpha_i+\beta_i\ne1
 \qquad(i\notin T),
\end{equation}
where outside $T$ the symbols are the integers $0,1$.
\end{definition}

\begin{theorem}[Prefort gluing]\label{thm:prefort-gluing}
Let $H,K$ be $q$-boundaried graphs with interior initial sets $S_H,S_K$, and
let
\[
 E=E_B(H)\cup E_B(K).
\]
There is a fort of $H\oplus_qK$ disjoint from
$S_H\cup S_K$ with boundary trace $T$ if and only if there exist
$E$-compatible atoms
\[
 \alpha\in\mathcal A_H(S_H),\qquad
 \beta\in\mathcal A_K(S_K)
\]
with common trace $T$.
\end{theorem}

\begin{proof}
Suppose first that $F$ is a global fort.  Its restrictions
$F_H=F\cap V(H)$ and $F_K=F\cap V(K)$ are local preforts and agree on the
boundary trace.  At a boundary vertex outside the trace, the total number of
fort neighbours is
\[
 c_i(F_H)+c_i(F_K)+d_E(i,T),
\]
and this number is not one.  If a local count is zero or one, take the forced
atomic value.  If it is at least two, choose a binary refinement so that
\eqref{eq:atomic-compatibility} holds; this is always possible because the
actual count is already at least two.  Compatible refinements exist.

Conversely, choose local preforts and atomic refinements producing compatible
$\alpha,\beta$.  Their union has the correct fort condition at every
nonboundary vertex by the definition of prefort.  The traces agree, so the
union is well defined at $B$.  At $i\notin T$, if both local counts are zero or
one, equation \eqref{eq:atomic-compatibility} is exactly the global fort
condition.  If either local count is at least two, the global count is at
least two automatically.  Hence the union is a global fort.
\end{proof}

For $A,D\in\mathfrak J_q$ let
\[
 \mathcal T_E(A,D)
 =\{T:\text{an $E$-compatible pair in $A\times D$ has trace }T\}.
\]
If $X\subseteq B$ is the initially blue part of the boundary, a global fort
disjoint from the initial set can only have trace disjoint from $X$.  We obtain
the separator kernel
\begin{equation}\label{eq:separator-kernel}
 K_E(A,D;x)
 =\sum_{\substack{X\subseteq B\\
      X\cap T\ne\varnothing\ \forall T\in\mathcal T_E(A,D)}}x^{|X|}.
\end{equation}
If $\varnothing\in\mathcal T_E(A,D)$ the sum is zero; if
$\mathcal T_E(A,D)=\varnothing$, it is $(1+x)^q$.

\begin{example}[A complete two-terminal calculation]\label{ex:q2-separator}
Let $q=2$, write the independent boundary as $B=\{a,b\}$, and consider the
two decorated pieces in \cref{fig:q2-worked}.  In the left piece $H$, the
interior vertices are $v,c$, the edges are
\[
 av,\ bv,\ ac,\ bc,
\]
and $S_H=\{c\}$.  Since the selected vertex $c$ lies outside every prefort
disjoint from $S_H$, its two neighbours $a,b$ must occur together or not at
all.  The nonempty preforts are therefore
\[
 \{v\},\qquad \{a,b\},\qquad \{a,b,v\}.
\]
They supply, respectively, the atoms $11,\mathsf I\mathsf I,$ and
$\mathsf I\mathsf I$.  Thus
\[
 \mathcal A_H(S_H)=\{11,\mathsf I\mathsf I\},
\]
a two-element chain.  In the right piece $K$, the only interior vertex $w$ is
selected and adjacent to both boundary vertices.  Its prefort condition
excludes $\{a\}$ and $\{b\}$, so
\[
 \mathcal A_K(S_K)=\{\mathsf I\mathsf I\},\qquad S_K=\{w\}.
\]

For an independent boundary, two atoms are compatible exactly when they are
equal.  Hence the only compatible pair is
$\mathsf I\mathsf I/\mathsf I\mathsf I$, with trace $\{a,b\}$.  A selected
boundary set must meet this trace, and therefore
\[
 K_{\varnothing}(\mathcal A_H(S_H),\mathcal A_K(S_K);x)
 =2x+x^2.
\]
Equivalently, the chain formula below gives the same value directly:
$\mathcal A_H(S_H)\cap\mathcal A_K(S_K)=\{\mathsf I\mathsf I\}$, so the
empty-chain term contributes $(1+x)^2$ and the singleton-chain term contributes
$-1$.  Thus $(1+x)^2-1=2x+x^2$.  The fixed interior choice
$S_H\cup S_K=\{c,w\}$ consequently contributes
\[
 x^2(2x+x^2)=2x^3+x^4
\]
to the polynomial of the glued graph.  Directly, one must initially select at
least one of $a,b$; either choice starts the forcing process, while selecting
neither leaves both boundary vertices white and no force is available.
\end{example}

\begin{figure}[ht]
\centering
\begin{tikzpicture}[every node/.style={font=\small},
 boundary/.style={circle,draw,double,minimum size=7mm,inner sep=0pt},
 plain/.style={circle,draw,minimum size=7mm,inner sep=0pt},
 selected/.style={circle,draw,fill=blue!18,minimum size=7mm,inner sep=0pt}]
 \node[boundary] (a1) at (0,1.1) {$a$};
 \node[boundary] (b1) at (3,1.1) {$b$};
 \node[plain] (v) at (1.5,2.0) {$v$};
 \node[selected] (c) at (1.5,0.2) {$c$};
 \draw (a1)--(v)--(b1)--(c)--(a1);
 \node at (1.5,-0.55) {$H,\ S_H=\{c\}$};

 \node[boundary] (a2) at (5.0,1.1) {$a$};
 \node[boundary] (b2) at (8.0,1.1) {$b$};
 \node[selected] (w) at (6.5,1.1) {$w$};
 \draw (a2)--(w)--(b2);
 \node at (6.5,-0.55) {$K,\ S_K=\{w\}$};
\end{tikzpicture}
\caption{Worked $q=2$ separator example.  Double circles are boundary
vertices; shaded vertices are prescribed members of the interior initial set.
The left atom family is the chain $\{11,\mathsf I\mathsf I\}$ and the right
family is $\{\mathsf I\mathsf I\}$.}
\label{fig:q2-worked}
\end{figure}

\begin{theorem}[Separator splitting formula]\label{thm:separator-splitting}
Let $H,K$ be $q$-boundaried graphs, write
$E=E_B(H)\cup E_B(K)$, and for $A\in\mathfrak J_q$ define
\[
 R_H(A;x)=
 \sum_{\substack{S\subseteq V(H)\setminus B\\
       0^q\notin\mathcal A_H(S),\ \mathcal A_H(S)=A}}x^{|S|},
\]
with $R_K$ defined analogously.  Then
\begin{equation}\label{eq:general-splitting}
 \boxed{
 \Zpoly(H\oplus_qK;x)
 =\sum_{A,D\in\mathfrak J_q}
   R_H(A;x)K_E(A,D;x)R_K(D;x).
 }
\end{equation}
\end{theorem}

\begin{proof}
Partition an initial set of $H\oplus_qK$ uniquely into its interior parts
$S_H,S_K$ and its selected boundary set $X$.  Fatal interior states cannot
succeed by \cref{lem:fatal-atom}.  For a nonfatal pair, the fort
characterization and \cref{thm:prefort-gluing} say that the total initial set
is zero forcing exactly when $X$ intersects every trace in
$\mathcal T_E(\mathcal A_H(S_H),\mathcal A_K(S_K))$.  Summing the weights
$x^{|S_H|+|S_K|+|X|}$ gives \eqref{eq:general-splitting}.
\end{proof}

We obtain a splitting formula for every finite separating set.  The rest of
this section determines the linear information actually needed by this finite
but potentially large kernel.

\section{Chain compression and exact connection ranks}\label{sec:separator-rank}

Let $\mathcal C_q$ be the set of chains of $P_q$, including the empty chain,
and put
\[
  c_q=|\mathcal C_q|.
\]
The key cancellation is independent of zero forcing.

\begin{lemma}[Chain disjointness identity]\label{lem:chain-disjointness}
If $A,D\in\mathfrak J_q$, then
\begin{equation}\label{eq:chain-disjointness}
 \mathbf1[A\cap D=\varnothing]
 =\sum_{\substack{C\in\mathcal C_q\\C\subseteq A\cap D}}(-1)^{|C|}.
\end{equation}
\end{lemma}

\begin{proof}
Ordinary inclusion--exclusion gives
\[
 \mathbf1[A\cap D=\varnothing]
 =\sum_{X\subseteq A\cap D}(-1)^{|X|}.
\]
For a finite nonempty join-semilattice $J$, call $j\in J$
\emph{join-irreducible} if there do not exist $u,v<j$ with $j=u\vee v$.
Let $I(J)$ be the set of join-irreducible elements.  Two elementary finite
semilattice facts will be used explicitly.

First, $I(J)$ generates $J$.  Indeed, proceed by induction upward from the
minimal elements.  A minimal element is join-irreducible.  If $j$ is not
join-irreducible, write $j=u\vee v$ with $u,v<j$ and apply the induction
hypothesis to $u$ and $v$.  Second, every generating set $X$ for $J$ must
contain $I(J)$.  If $j\in I(J)$ is a join of finitely many elements of $X$,
repeated use of join-irreducibility forces $j$ to equal one of those
generators.

Now group the subsets $X\subseteq A\cap D$ according to the join-closed set
$J=\langle X\rangle_\vee$ that they generate.  For nonempty $J$, the preceding
facts show that the generating subsets are exactly
\[
 I(J)\cup Y,\qquad Y\subseteq J\setminus I(J).
\]
Their total inclusion--exclusion coefficient is therefore
\[
 \sum_{Y\subseteq J\setminus I(J)}(-1)^{|I(J)|+|Y|}
 =(-1)^{|I(J)|}(1-1)^{|J|-|I(J)|}.
\]
This is nonzero precisely when every element of $J$ is join-irreducible.
That happens precisely when $J$ is a chain.  If $a,b\in J$ are incomparable,
then $a\vee b$ is strictly larger than both and is therefore reducible.
Conversely, in a chain the join of two elements is their maximum, so no element
can be the join of two strictly smaller elements.  A nonempty chain $C$
contributes $(-1)^{|C|}$.

Finally, $X=\varnothing$ generates the empty join-closed set; its coefficient
is $1$, which is exactly the contribution of the empty chain in
\eqref{eq:chain-disjointness}.  This proves the identity including the empty
case.
\end{proof}

Define the chain-incidence matrix
\[
 Z[A,C]=\mathbf1[C\subseteq A],
 \qquad A\in\mathfrak J_q,
 \quad C\in\mathcal C_q.
\]
When its rows are restricted to chain states, $Z$ is the zeta matrix of the
finite poset $\mathcal C_q$ ordered by inclusion.  After a linear extension it
is triangular with diagonal one, hence invertible.

The next lemma is what allows chain compression to survive arbitrary boundary
edges.

\begin{lemma}[Compatibility neighbourhoods are join closed]
\label{lem:compatibility-neighborhood}
Fix a boundary-edge graph $E$, a selected boundary set $X$, and
$D\in\mathfrak J_q$.  Let $N_{E,X}(D)$ be the set of atoms $\alpha\in P_q$
whose trace is disjoint from $X$ and for which some $\beta\in D$ is
$E$-compatible with $\alpha$.  Then $N_{E,X}(D)$ is join closed.
\end{lemma}

\begin{proof}
Take $\alpha_1,\alpha_2\in N_{E,X}(D)$ with compatible witnesses
$\beta_1,\beta_2\in D$.  Since $D$ is join closed,
$\beta=\beta_1\vee\beta_2$ belongs to $D$; put
$\alpha=\alpha_1\vee\alpha_2$.  Their common trace is the union of the two
original traces and is still disjoint from $X$.

Suppose compatibility failed at $i$ outside that union.  If the boundary-edge
contribution $d_E(i,T)$ were one and both joined bits zero, then one of the
original traces already contributed that unique boundary neighbour while the
corresponding original bits were zero, contradicting compatibility of that
pair.  If $d_E(i,T)=0$ and the two joined bits summed to one, then one joined
bit, say on the $\alpha$ side, is one while the other is zero.  Some original
$\alpha_j$ therefore has bit one, every $\beta_j$ has bit zero, and no original
trace has an $E$-neighbour at $i$; again the $j$th pair would violate
\eqref{eq:atomic-compatibility}.  Contributions at least two are automatically
safe.  The atoms $\alpha$ and $\beta$ are compatible.
\end{proof}

For fixed $E$ and $X$, the success indicator is therefore
\[
 k_{E,X}(A,D)
 =\mathbf1[A\cap N_{E,X}(D)=\varnothing].
\]
By \cref{lem:chain-disjointness}, its dependence on $A$ lies in the column
space of $Z$.  By symmetry the same is true in the $D$ coordinate.  Summing
$x^{|X|}k_{E,X}$ over $X$ proves the following.

\begin{theorem}[Chain compression]\label{thm:chain-compression}
For every boundary-edge graph $E$, the matrix
\[
 \bigl(K_E(A,D;x)\bigr)_{A,D\in\mathfrak J_q}
\]
over $\mathbb Q(x)$ has the same rank as its submatrix obtained by restricting
both rows and columns to chain states $C,D\in\mathcal C_q$.  In particular its
rank is at most $c_q$.
\end{theorem}

\begin{proof}
The preceding discussion gives a factorization $K_E=ZH_E$ for some matrix
$H_E$.  Since the restriction of $Z$ to chain rows is invertible, the chain
rows span the full row space.  Applying the same argument to the symmetric
matrix $K_E$ shows that the chain columns span the full column space.  Hence
the chain--chain submatrix has the same rank.
\end{proof}

\subsection{Independent boundaries and chain enumeration}

When $E=\varnothing$, two atoms are compatible exactly when they are equal.
Indeed they must have the same trace, and outside the trace
$\alpha_i+\beta_i\ne1$ is equivalent to $\alpha_i=\beta_i$ for binary bits.
Inclusion--exclusion over the common atoms gives the more explicit
factorization
\begin{equation}\label{eq:independent-chain-factor}
 K_{\varnothing}(A,D;x)
 =\sum_{\substack{C\in\mathcal C_q\\C\subseteq A\cap D}}
 (-1)^{|C|}(1+x)^{q-|T(C)|},
\end{equation}
where $T(C)=\bigcup_{\alpha\in C}T(\alpha)$.  Thus
\[
 K_{\varnothing}=Z\,W(x)\,Z^T,
 \qquad
 W_{C,C}(x)=(-1)^{|C|}(1+x)^{q-|T(C)|}.
\]
Every diagonal entry is nonzero in $\mathbb Q(x)$ and the chain rows of $Z$
are invertible.  Therefore
\begin{equation}\label{eq:independent-universal-rank}
 \operatorname{rank}_{\mathbb Q(x)}K_{\varnothing}=c_q.
\end{equation}

The numbers $c_q$ are explicit.

\begin{proposition}[Chain count]\label{prop:chain-count}
For $q\ge0$,
\begin{equation}\label{eq:cq-formula}
 \boxed{
 c_q=\frac12\left[1+\sum_{k=1}^{2q+1}\sum_{j=1}^k
 (-1)^{k-j}\binom{k-1}{j-1}\binom{j+2}{2}^{q}\right].
 }
\end{equation}
Moreover
\[
 \log c_q=2q\log q+O(q).
\]
\end{proposition}

\begin{proof}
A weak $j$-term chain in $C_3^q$ is chosen independently in every coordinate.
The number of nondecreasing $j$-tuples from a three-element chain is
$\binom{j+2}{2}$, so there are $\binom{j+2}{2}^q$ weak chains.  If $s_q(k)$ is
the number of strict $k$-element chains, then a weak $j$-chain with exactly
$k$ distinct terms is obtained by choosing the $k-1$ change positions, giving
\[
 \binom{j+2}{2}^q
 =\sum_{k=1}^j\binom{j-1}{k-1}s_q(k).
\]
Binomial inversion yields
\[
 s_q(k)=\sum_{j=1}^k(-1)^{k-j}\binom{k-1}{j-1}\binom{j+2}{2}^q.
\]
A strict chain in $C_3^q$ has at most $2q+1$ elements.  Since $C_3^q$ has a
unique minimum $0^q$, adding or deleting that minimum bijects chains avoiding
it with chains containing it.  Hence the number of chains in
$P_q=C_3^q\setminus\{0^q\}$ is one half of the total number of chains in
$C_3^q$, proving \eqref{eq:cq-formula}.

For the asymptotics, maximal chains are obtained by ordering the $2q$ labeled
coordinate increments, subject only to the first increment of each coordinate
preceding its second.  There are $(2q)!/2^q$ such chains.  Conversely, encode
an arbitrary chain by assigning each of the $2q$ coordinate increments a time
in $\{0,1,\ldots,2q+1\}$, where the extreme values mean that the increment
occurs before the first or after the last displayed element.  This gives the
upper bound $(2q+2)^{2q}$.  Stirling's formula now gives the stated logarithmic
order.
\end{proof}

The first values are
\[
\begin{array}{c|rrrrrrrr}
q&0&1&2&3&4&5&6&7\\
\hline
c_q&1&4&52&1636&95668&8977444&1234728052&234059341156.
\end{array}
\]

\subsection{Realizing every chain state}

The chain compression above is an upper-bound statement unless its chain modes
can be produced by actual graph pieces.  We next give an explicit realization.
A \emph{decorated graph} means a graph together with a prescribed interior
initial set $S$; selected constraint vertices below belong to $S$ and therefore
cannot lie in a prefort disjoint from $S$.

\begin{lemma}[Equality and implication gadgets]\label{lem:boolean-gadgets}
Let $u,v$ be unselected nonboundary vertices.
\begin{enumerate}[label=\textnormal{(\alph*)}]
\item A selected vertex adjacent exactly to $u,v$ enforces
$u\in F\Longleftrightarrow v\in F$ in every prefort $F$ disjoint from the
selected set.
\item After adding an unselected duplicate $v'$ equality-tied to $v$, a
selected vertex adjacent to $u,v,v'$ enforces
$u\in F\Longrightarrow v\in F$.
\end{enumerate}
\end{lemma}

\begin{proof}
A selected constraint vertex lies outside $F$, so its prefort condition says
that the number of its neighbours in $F$ is not one.  For degree two this
allows only $00$ and $11$, proving equality.  In part (b), the duplicate has
the same membership as $v$, so the three-neighbour counts for $(u,v)$ are
$0,2,1,3$ on $00,01,10,11$, respectively.  Exactly the assignment $10$ is
forbidden.
\end{proof}

\begin{theorem}[Chain realization]\label{thm:chain-realization}
For every chain $C\in\mathcal C_q$ and every desired boundary-edge set
$E\subseteq\binom B2$, there exists a finite simple $q$-boundaried graph $H_C$
and an interior initial set $S_C$ such that
\[
 E_B(H_C)=E,
 \qquad
 0^q\notin\mathcal A_{H_C}(S_C),
 \qquad
 \mathcal A_{H_C}(S_C)=C.
\]
\end{theorem}

\begin{proof}
We give the construction with a formal vertex and edge description.  All
vertices called \emph{selected} below belong to $S_C$; all other auxiliary
vertices are unselected.

\paragraph{The empty chain.}
For $C=\varnothing$, take the boundary $B$ and, for every $i$, add a selected
leaf $p_i$ adjacent only to $b_i$.  If a prefort disjoint from $S_C$ contained
$b_i$, then the selected vertex $p_i$, which lies outside that prefort, would
have exactly one prefort neighbour.  Hence every boundary vertex is excluded.
There are no other unselected vertices, so no nonempty prefort exists and the
atomic family is empty.  The prescribed boundary edges $E$ may then be added.

\paragraph{The nonempty chain: master variables.}
Write
\[
 C=\{\alpha^1<\alpha^2<\cdots<\alpha^k\},\qquad k\ge1.
\]
Introduce unselected master vertices $t_1,\ldots,t_k$.  For every
$1\le j<k$, introduce an unselected duplicate $d_j$ and selected vertices
$e_j,p_j$, with edges
\[
 e_jt_j,\ e_jd_j,
 \qquad
 p_jt_{j+1},\ p_jt_j,\ p_jd_j.
\]
The degree-two selected vertex $e_j$ enforces
$d_j\in F\Longleftrightarrow t_j\in F$, and then the degree-three selected
vertex $p_j$ enforces
\[
 t_{j+1}\in F\Longrightarrow t_j\in F
\]
by \cref{lem:boolean-gadgets}.  Hence the possible membership vectors of the
master vertices are exactly the prefixes
\[
 00\cdots0,\ 10\cdots0,\ 110\cdots0,\ldots,\ 11\cdots1.
\]

\paragraph{Encoding one boundary coordinate.}
Fix $i\in B$.  Because $C$ is a chain, the coordinate sequence
$\alpha_i^1,\ldots,\alpha_i^k$ is nondecreasing in
$0<1<\mathsf I$.  If the value $1$ occurs, let
\[
 a(i)=\min\{j:\alpha_i^j=1\}.
\]
Introduce an unselected marker $m_i$, join it to $b_i$, and introduce a
selected equality vertex $f_i$ with edges $f_im_i,f_it_{a(i)}$.  Thus
$m_i\in F$ exactly from prefix $a(i)$ onward.  If the value $1$ never occurs,
no marker is introduced.

If the value $\mathsf I$ occurs, let
\[
 \iota(i)=\min\{j:\alpha_i^j=\mathsf I\}.
\]
Introduce a selected equality vertex $g_i$ with edges
$g_ib_i,g_it_{\iota(i)}$; hence $b_i$ enters the prefort exactly from prefix
$\iota(i)$ onward.  If $\mathsf I$ never occurs, instead introduce a selected
pin $h_i$ adjacent only to $b_i$, which forces $b_i$ to remain outside every
prefort.  A direct jump $0\to\mathsf I$ requires no marker.

Let $S_C$ be the set of all selected vertices
\[
 \{e_j,p_j\}_{j<k}\cup\{f_i,g_i,h_i:\text{created above}\},
\]
and let $H_C$ contain precisely the boundary, master, duplicate, marker, and
selected constraint vertices just described, together with the listed edges.
Finally add the prescribed edge set $E$ inside the boundary.  These boundary
edges do not enter any nonboundary prefort condition.

\paragraph{No unintended preforts.}
Let $F$ be a prefort disjoint from $S_C$.  The equality and implication
constraints force the master membership set to be a prefix
$\{t_1,\ldots,t_s\}$ for some $0\le s\le k$, and all duplicates, markers, and
unpinned boundary vertices are then determined by equality.  If $s=0$, every
duplicate and marker is outside $F$, every equality-tied boundary vertex is
outside $F$, and every remaining boundary vertex is pinned out; hence $F$ is
empty, contrary to the definition of prefort.  Thus $1\le s\le k$.

Conversely, each nonzero prefix gives a prefort.  Every selected equality or
implication vertex sees one of the allowed neighbour counts $0,2,$ or $3$ by
construction.  An unselected master vertex or duplicate that lies outside the
prefort is adjacent only to selected constraint vertices and therefore sees
zero prefort neighbours.  The only additional check concerns a marker $m_i$
outside the prefort.  This means $s<a(i)$.  Monotonicity of the coordinate
sequence implies that $b_i$ has not yet reached $\mathsf I$, so $b_i$ is also
outside the prefort.  The marker therefore sees zero prefort neighbours (its
other neighbour is the selected equality vertex $f_i$).  Boundary vertices
need not satisfy a local fort condition by definition of prefort.  Hence no
unselected auxiliary vertex creates an unintended restriction.

For prefix $s$, coordinate $i$ is $0$ before $a(i)$, is $1$ from $a(i)$ until
$\iota(i)-1$ when these stages exist, and is $\mathsf I$ from $\iota(i)$
onward.  The resulting atomic word is exactly $\alpha^s$.  Each boundary
vertex has at most one interior prefort neighbour, so no $\ge2$ count creates
extra wildcard refinements.  Therefore
\[
 \mathcal A_{H_C}(S_C)=\{\alpha^1,\ldots,\alpha^k\}=C.
\]
Since $C\subseteq P_q$, the fatal atom $0^q$ does not occur.  Adding the edges
of $E$ inside $B$ changes neither the prefort conditions away from $B$ nor the
interior counts $c_i(F)$, completing the construction.
\end{proof}

\begin{figure}[ht]
\centering
\begin{tikzpicture}[every node/.style={font=\scriptsize},
 var/.style={circle,draw,minimum size=6mm,inner sep=0pt},
 boundary/.style={circle,draw,double,minimum size=6mm,inner sep=0pt},
 gadget/.style={draw,rounded corners,fill=blue!10,minimum height=5mm,inner sep=2pt}]
 \node[var] (t1) at (0,0) {$t_1$};
 \node[var] (t2) at (2,0) {$t_2$};
 \node[gadget] (imp) at (1,0.75) {$t_2\Rightarrow t_1$};
 \draw (t1)--(imp)--(t2);
 \node[var] (m1) at (-0.5,2) {$m_1$};
 \node[var] (m2) at (2.5,2) {$m_2$};
 \node[gadget] (eqm1) at (-0.25,1) {$=$};
 \node[gadget] (eqm2) at (2.25,1) {$=$};
 \draw (t1)--(eqm1)--(m1);
 \draw (t1)--(eqm2)--(m2);
 \node[boundary] (b1) at (-0.5,3.2) {$b_1$};
 \node[boundary] (b2) at (2.5,3.2) {$b_2$};
 \draw (m1)--(b1); \draw (m2)--(b2);
 \node[gadget] (eqb1) at (0.25,3.2) {$=$};
 \node[gadget] (eqb2) at (1.75,3.2) {$=$};
 \draw (b1)--(eqb1)--(t2);
 \draw (b2)--(eqb2)--(t2);
 \node[align=left] at (5.2,1.6) {schematic canonical realization\\of $C=\{11,\mathsf I\mathsf I\}$;\\boxes denote the selected gadgets\\of \cref{lem:boolean-gadgets}};
\end{tikzpicture}
\caption{Schematic of the chain-realization construction for a representative
$q=2$ chain.  The implication box abbreviates the duplicate-plus-constraint
gadget specified formally in the proof; equality boxes are selected
degree-two constraint vertices.}
\label{fig:chain-realization}
\end{figure}

\subsection{Selection projectors}

To turn a decorated realization into a row of the ordinary connection matrix,
we must isolate the prescribed initial set.  The rank-four one-terminal
calculus provides explicit projectors.  We work with formal
$\mathbb Q(x)$-linear combinations of rooted graphs, usually called quantum
graphs in the connection-matrix literature.

Let $K_1,K_2,P_3^{\rm end},P_3^{\rm mid}$ denote the four rooted contexts used
in \cref{thm:connection-rank}.  Define
\begin{align}
 Q_{\rm in}
 &= -\frac1{x+1}K_1
    -\frac1{x(x+1)}K_2
    +\frac1{x(x+1)}P_3^{\rm end},
 \label{eq:projector-in}\\
 Q_{\rm out}
 &= \frac{x+2}{x+1}K_1
    +\frac1{x(x+1)}K_2
    -\frac1{x(x+1)}P_3^{\rm end}.
 \label{eq:projector-out}
\end{align}

\begin{lemma}[Vertex-selection projectors]\label{lem:selection-projectors}
For every rooted graph $(G,v)$,
\begin{align*}
 \Zpoly(G\vee Q_{\rm in};x)
 &=\sum_{\substack{S\text{ zero forcing in }G\\v\in S}}x^{|S|},\\
 \Zpoly(G\vee Q_{\rm out};x)
 &=\sum_{\substack{S\text{ zero forcing in }G\\v\notin S}}x^{|S|}.
\end{align*}
The identities remain valid when $G$ contains arbitrary additional labeled
boundary vertices and arbitrary external attachments away from the interior of
the projector.
\end{lemma}

\begin{proof}
Write $r(G)=(P_0,P_1,A_0,A_1)^T$.  The rooted profiles used in the projectors
are
\[
 r(K_1)=\begin{pmatrix}1\\0\\0\\0\end{pmatrix},\qquad
 r(K_2)=\begin{pmatrix}0\\x\\1\\0\end{pmatrix},\qquad
 r(P_3^{\rm end})=\begin{pmatrix}x\\x+x^2\\1\\0\end{pmatrix}.
\]
Consequently the formal response vectors of the two quantum graphs simplify to
\[
 r(Q_{\rm in})=
 \begin{pmatrix}0\\x/(x+1)\\0\\0\end{pmatrix},
 \qquad
 r(Q_{\rm out})=
 \begin{pmatrix}1\\-x/(x+1)\\0\\0\end{pmatrix}.
\]
Multiplication by the connection matrix $B(x)$ of
\cref{thm:connection-rank} gives
\[
 B(x)r(Q_{\rm in})=
 \begin{pmatrix}x\\x\\x\\x\end{pmatrix},
 \qquad
 B(x)r(Q_{\rm out})=
 \begin{pmatrix}0\\1\\0\\1\end{pmatrix}.
\]
Hence
\[
 r(G)^TB(x)r(Q_{\rm in})=x(P_0+P_1+A_0+A_1),
 \qquad
 r(G)^TB(x)r(Q_{\rm out})=P_1+A_1,
\]
which are exactly the contributions from initial sets containing and omitting
$v$, respectively.  All calculations take place over $\mathbb Q(x)$, so the
denominators in the projectors are legitimate; specializations such as
$x=0,-1$ are separate questions.

The argument is coefficientwise in any additional vertex-set variables on
$G-v$.  Thus extra labeled boundary vertices or external attachments can be
held fixed as part of the rooted context, and successive projectors at
distinct vertices act as commuting membership filters.
\end{proof}

Applying these projectors successively at distinct vertices filters the
multivariate zero-forcing generating function coefficientwise, so they commute
as membership filters.

\begin{corollary}[Pinning a prescribed interior set]\label{cor:pinning}
Let $H$ be a $q$-boundaried graph, $U=V(H)\setminus B$, and
$S\subseteq U$.  There is a $q$-boundaried quantum graph $\Pi_{H,S}$ which, in every
external $q$-boundaried context, retains exactly the zero-forcing sets whose
intersection with $U$ is $S$.  In particular, if $(H,S)$ realizes a chain
state $C$ and $(K,T)$ realizes a chain state $D$, then
\[
 \Zpoly(\Pi_{H,S}\oplus_q\Pi_{K,T};x)
 =x^{|S|+|T|}K_{E_B(H)\cup E_B(K)}(C,D;x).
\]
\end{corollary}

\begin{proof}
Attach $Q_{\rm in}$ at every vertex of $S$ and $Q_{\rm out}$ at every vertex
of $U\setminus S$.  Repeated application of
\cref{lem:selection-projectors} retains exactly the desired membership
pattern, and the retained interior vertices contribute $x^{|S|}$.  Applying
the same construction independently to $(K,T)$ and then using the separator
splitting theorem leaves only the prescribed pair of local states, giving the
displayed pairing identity.
\end{proof}

\subsection{The master connection-rank formula}

For a fixed boundary-edge graph $E$, let
\[
  M_E(x)=\bigl(K_E(C,D;x)\bigr)_{C,D\in\mathcal C_q}
  \in\mathbb Q(x)^{c_q\times c_q}
\]
be the chain kernel.  Put $m=\binom q2$.  The full finite
\emph{chain-edge connection kernel}
$\mathbb M_q\in\mathbb Q(x)^{2^m c_q\times 2^m c_q}$ is indexed by pairs
$(F,C)$ with $F\subseteq\binom B2$ and $C\in\mathcal C_q$ and has block
\[
 \mathbb M_q[(F,C),(G,D)]=M_{F\cup G}(x)[C,D].
\]

\begin{theorem}[Exact finite connection model]\label{thm:finite-connection-model}
For every $q\ge0$,
\[
 \boxed{
 r_q=\operatorname{rank}_{\mathbb Q(x)}\mathbb M_q.
 }
\]
The infinite connection matrix of the zero-forcing polynomial has exactly
the same rank as the explicit finite chain-edge kernel.
\end{theorem}

\begin{proof}
Put $m=\binom q2$ and $j_q=|\mathfrak J_q|$.  We prove the two inequalities
separately and keep the finite and infinite column spaces distinct throughout.

\emph{Upper bound.}
For a fixed glued boundary-edge graph $E$, let
\[
 U_E=\bigl(K_E(A,D;x)\bigr)_{A,D\in\mathfrak J_q}
 \in\mathbb Q(x)^{j_q\times j_q}.
\]
Let $Z\in\{0,1\}^{j_q\times c_q}$ be the join-state/chain incidence matrix
$Z[A,C]=\mathbf1[C\subseteq A]$, and let $Z_0$ be its square submatrix on
chain rows.  As observed above, $Z_0$ is the zeta matrix of chains ordered by
inclusion and is invertible; in particular $Z$ has full column rank.

The proof of \cref{thm:chain-compression} shows that the column and row spaces
of $U_E$ are contained in the column and row spaces generated by $Z$.  Choose
any left inverse $L$ of $Z$ and set
\[
 N_E=L U_E L^T\in\mathbb Q(x)^{c_q\times c_q}.
\]
Since $ZL$ acts as the identity on the column space of $U_E$ and
$L^TZ^T$ acts as the identity on its row space,
\begin{equation}\label{eq:UE-ZNE}
 U_E=ZN_EZ^T.
\end{equation}
Restricting \eqref{eq:UE-ZNE} to chain rows and columns gives
\begin{equation}\label{eq:ME-Z0NE}
 M_E=Z_0N_EZ_0^T.
\end{equation}
The matrices $N_E$ and $M_E$ are congruent through the invertible matrix $Z_0$.

Let $\mathbb U_q$ be the universal finite matrix indexed by local states
$(F,A)$ with $F\subseteq\binom B2$ and $A\in\mathfrak J_q$, whose
$(F,G)$ block is $U_{F\cup G}$.  Let $\mathbb N_q$ be the analogous block
matrix with blocks $N_{F\cup G}$.  Equations \eqref{eq:UE-ZNE} and
\eqref{eq:ME-Z0NE} give the explicit factorizations
\begin{align}
 \mathbb U_q
 &=(I_{2^m}\otimes Z)\,\mathbb N_q\,(I_{2^m}\otimes Z)^T,
 \label{eq:universal-state-factor}\\
 \mathbb M_q
 &=(I_{2^m}\otimes Z_0)\,\mathbb N_q\,(I_{2^m}\otimes Z_0)^T.
 \label{eq:chain-state-factor}
\end{align}
Because $Z$ has full column rank and $Z_0$ is invertible,
\[
 \operatorname{rank}\mathbb U_q
 =\operatorname{rank}\mathbb N_q
 =\operatorname{rank}\mathbb M_q.
\]

For each ordinary $q$-boundaried graph $H$, form its finite response vector
$\rho_H$ indexed by $(F,A)$: it is zero unless $F=E_B(H)$, and at
$(E_B(H),A)$ it equals the polynomial $R_H(A;x)$ from
\cref{thm:separator-splitting}.  The splitting formula is precisely
\[
 C_q[H,K]=\rho_H^T\mathbb U_q\rho_K.
\]
Hence the infinite connection matrix factors through $\mathbb U_q$, and
\[
 r_q\le\operatorname{rank}\mathbb U_q
 =\operatorname{rank}\mathbb M_q.
\]

\emph{Lower bound.}
Index the chain-edge states by $i=(F,C)$.  By
\cref{thm:chain-realization}, choose for each $i$ a decorated realization
$(H_i,S_i)$ with boundary-edge set $F$ and chain state $C$, and put
\[
 Q_i=\Pi_{H_i,S_i}.
\]
Apply the pinning construction on \emph{both} sides.  If
$i=(F,C)$ and $j=(G,D)$, \cref{cor:pinning} gives
\begin{equation}\label{eq:pinned-pairing}
 \Zpoly(Q_i\oplus_qQ_j;x)
 =x^{|S_i|+|S_j|}K_{F\cup G}(C,D;x).
\end{equation}
Therefore the finite pairing matrix of the quantum graphs $Q_i$ is
\begin{equation}\label{eq:DMD-pairing}
 \bigl(\Zpoly(Q_i\oplus_qQ_j;x)\bigr)_{i,j}
 =D\,\mathbb M_q\,D,
 \qquad D_{ii}=x^{|S_i|}.
\end{equation}
The diagonal matrix $D$ is invertible over $\mathbb Q(x)$.

Each $Q_i$ is a finite $\mathbb Q(x)$-linear combination of ordinary
$q$-boundaried graphs.  Taking the finite union of all ordinary graphs that
occur in these combinations, let $P$ be their coefficient matrix and let
$C_q^{\rm fin}$ be the corresponding finite submatrix of the ordinary
connection matrix.  Bilinearity of gluing gives
\[
 D\mathbb M_qD=P\,C_q^{\rm fin}P^T.
\]
Consequently
\[
 \operatorname{rank}\mathbb M_q
 =\operatorname{rank}(D\mathbb M_qD)
 \le \operatorname{rank}C_q^{\rm fin}
 \le r_q.
\]
Together with the upper bound, this proves the theorem.
\end{proof}

The Boolean transform makes this rank additive over edge sectors.  For
$R\subseteq\binom B2$ define
\begin{equation}\label{eq:edge-mobius}
 \widehat M_R(x)
 =\sum_{E\supseteq R}(-1)^{|E\setminus R|}M_E(x).
\end{equation}

\begin{theorem}[Master rank formula]\label{thm:master-rank}
For every $q\ge0$,
\begin{equation}\label{eq:master-rank}
 \boxed{
 r_q=
 \sum_{R\subseteq\binom{[q]}2}
 \operatorname{rank}_{\mathbb Q(x)}\widehat M_R(x).
 }
\end{equation}
In particular, if $r_q^{\circ}$ denotes the connection rank of the
submatrix obtained by restricting \emph{both} local sides to $q$-boundaried
graphs whose boundary induces no edge, then
\begin{equation}\label{eq:independent-rank}
 \boxed{r_q^{\circ}=c_q.}
\end{equation}
Consequently
\begin{equation}\label{eq:rank-bounds}
 c_q\le r_q\le 2^{\binom q2}c_q.
\end{equation}
\end{theorem}

\begin{proof}
Let $L=\binom B2$ and index the subsets of $L$ in any linear extension of
inclusion.  Define the Boolean zeta matrix
\[
 Z_{\rm edge}[F,R]=\mathbf1[F\subseteq R],
 \qquad F,R\subseteq L.
\]
It is triangular with diagonal one and therefore invertible.  By the definition
\eqref{eq:edge-mobius}, Boolean M\"obius inversion gives
\[
 M_E=\sum_{R\supseteq E}\widehat M_R.
\]
Hence the $(F,G)$ block of $\mathbb M_q$ satisfies
\[
 M_{F\cup G}
 =\sum_{R\subseteq L}
   \mathbf1[F\subseteq R]\mathbf1[G\subseteq R]\widehat M_R.
\]
Equivalently,
\begin{equation}\label{eq:boolean-edge-congruence}
 \boxed{
 \mathbb M_q
 =(Z_{\rm edge}\otimes I_{c_q})
   \operatorname{diag}_{R\subseteq L}(\widehat M_R)
   (Z_{\rm edge}\otimes I_{c_q})^T.
 }
\end{equation}
Since the two zeta factors are invertible, rank is invariant under this
congruence and equals the sum of the ranks of the diagonal blocks.  Combining
with \cref{thm:finite-connection-model} proves \eqref{eq:master-rank}.

For a one-edge boundary $L=\{e\}$, \eqref{eq:boolean-edge-congruence} is the
concrete sanity check
\[
 \begin{pmatrix}
 M_{\varnothing}&M_e\\
 M_e&M_e
 \end{pmatrix}
 \sim
 \begin{pmatrix}
 M_{\varnothing}-M_e&0\\
 0&M_e
 \end{pmatrix}.
\]
For the restricted connection matrix in which the boundary is required to be
independent on each local side, only the single block $M_{\varnothing}$ is
present.  Its rank is $c_q$ by \eqref{eq:independent-universal-rank} together
with the realization-and-pinning lower bound.  Finally, every
$\widehat M_R$ is a $c_q\times c_q$ matrix, which gives
\eqref{eq:rank-bounds}.
\end{proof}

The quotient of the vector space of $q$-boundaried graphs by the radical of
$C_q$ is therefore a finite-dimensional commutative graph algebra of dimension
$r_q$ in the standard connection-matrix sense \cite{Lovasz2006}.  The case
$q=1$ is the explicit four-dimensional algebra of
\cref{thm:connection-rank}; higher $q$ need not admit such a small presentation,
but \cref{thm:master-rank} determines their dimensions by a finite
zero-forcing-specific kernel.

\subsection{The first two nontrivial ranks}

For $q=1$ there are no boundary edges and $c_1=4$, so the general theory
recovers
\[
 r_1=4.
\]
For $q=2$, the chain set has size $c_2=52$ and there is one possible boundary
edge $e$.  The Boolean transform has two sectors,
\[
 \widehat M_{\{e\}}=M_{\{e\}},
 \qquad
 \widehat M_{\varnothing}=M_{\varnothing}-M_{\{e\}}.
\]

\begin{proposition}[Machine-certified $q=2$ sector ranks]
\label{prop:q2-certified-ranks}
The exact finite matrices defined above satisfy
\begin{equation}\label{eq:q2-ranks}
 \operatorname{rank}_{\mathbb Q(x)}M_{\{e\}}=44,
 \qquad
 \operatorname{rank}_{\mathbb Q(x)}(M_{\varnothing}-M_{\{e\}})=16.
\end{equation}
\end{proposition}

\begin{proof}
This proposition is a finite exact computation, not an additional infinite
structural argument.  The supplementary certificate
\path{separator_rank_certificate.json} fixes the chain-state order,
canonical hashes of the integer coefficient matrices, the specialization
$x=2$, the prime $1{,}000{,}003$, and explicit nonsingular minor indices.
The independent verifier
\path{verify_separator_rank_certificate.py} reconstructs all matrices
from the definitions in this paper.  Exact rational row reduction of the
concatenated coefficient matrices gives upper ranks $44$ and $16$.  At $x=2$,
exact elimination over $\mathbb F_{1{,}000{,}003}$ verifies nonsingular minors
of orders $44$ and $16$, respectively, giving matching lower bounds because
specialization cannot increase generic rank.  No floating-point rank decision
is used.
\end{proof}

\begin{corollary}\label{cor:r2}
\begin{equation}\label{eq:r2}
 \boxed{r_2=60.}
\end{equation}
\end{corollary}

\begin{proof}
Apply \cref{thm:master-rank,prop:q2-certified-ranks}:
$r_2=44+16=60$.
\end{proof}


\section{Pendant-edge transfer}\label{sec:edge-extension}

The rooted product combines branches sharing a root.  To recurse along a tree, we also need to move the root through an edge.

Let $(H,u)$ be a rooted graph.  Form a new graph $E(H,u)$ by adjoining a new vertex $r$ adjacent only to $u$, and root the resulting graph at $r$.  We call this operation \emph{edge extension}.  Let
\[
  (P_0,P_1,A_0,A_1)
\]
be the profile of $(H,u)$, and let
\[
  (\widehat P_0,\widehat P_1,\widehat A_0,\widehat A_1)
\]
be the profile of $E(H,u)$.

\begin{theorem}[Edge-extension transfer]\label{thm:edge-extension}
The response polynomials satisfy
\begin{align}
  \widehat P_0&=A_1+xA,\label{eq:edgeP0}\\
  \widehat P_1&=P_1+xP,\label{eq:edgeP1}\\
  \widehat A_0&=P_0+A_0,\label{eq:edgeA0}\\
  \widehat A_1&=0.\label{eq:edgeA1}
\end{align}
\end{theorem}

\begin{proof}
Let $T\subseteq V(H)\setminus\{u\}$.  A set counted by the new profile either contains $u$ or does not.

Suppose first that $u$ is initially blue; the corresponding sets are $T\cup\{u\}$ and contribute a factor $x$.  With the new root $r$ blue and forbidden from forcing, the evolution inside $H$ is precisely the rooted process at $u$.  Hence the new branch is passive whenever $T$ is passive or active in $(H,u)$.  If $T$ is passive, then $u$ does not spend its force inside $H$; after $H$ has completed, $u$ can force an initially white $r$.  These sets contribute $xP$ to $\widehat P_1$.  If $T$ is active, then completion of $H$ consumes the unique force by $u$.  With $r$ initially white, the additional white neighbour $r$ prevents the required internal force, and the source indicator is zero.  These sets contribute $xA$ to $\widehat P_0$.

Now suppose that $u$ is initially white.  If $T$ lies in $P_1$, then vertices of $H-u$ force $u$ and complete $H$ without a force by $u$; the vertex $u$ may then force $r$.  Hence $P_1$ contributes to $\widehat P_1$.  If $T$ lies in $A_1$, then vertices of $H-u$ force $u$, but the subsequent completion of $H$ requires the unique force by $u$.  The extra white neighbour $r$ blocks that force when $r$ is initially white.  Hence these sets are passive with source indicator zero, contributing $A_1$ to $\widehat P_0$.

Finally, if $T$ lies in $P_0\cup A_0$, then $T$ does not force $u$.  With $r$ initially blue, the force $r\to u$ is both necessary and sufficient to place $H$ in the corresponding successful rooted state.  The new branch is therefore active and has source indicator zero.  This contributes $P_0+A_0$ to $\widehat A_0$.  An active state cannot source a pendant root, so $\widehat A_1=0$.
\end{proof}

It is useful to regard \cref{thm:edge-extension} as a linear operator on the four-component profile.  The rooted product is bilinear, while edge extension is linear; together they constitute a finite transfer calculus for trees.

\subsection{Elementary checks}\label{subsec:sanity}

The following profiles provide useful checks on the conventions.  The entries
are ordered as $(P_0,P_1,A_0,A_1)$.
\[
\begin{array}{c|c|c}
(H,r)&(P_0,P_1,A_0,A_1)&\Zpoly(H;x)\\
\hline
K_1&(1,0,0,0)&x\\
K_2\text{, root at an endpoint}&(0,x,1,0)&2x+x^2\\
P_3\text{, root at an endpoint}&(x,x+x^2,1,0)&2x+3x^2+x^3\\
P_3\text{, root at the middle}&(0,x^2,0,2x)&2x+3x^2+x^3\\
C_3&(2x,x^2,0,0)&3x^2+x^3\\
C_4&(2x+x^2,2x^2+x^3,0,0)&4x^2+4x^3+x^4
\end{array}
\]
For the star $K_{1,d}$ rooted at its centre, with $d\ge2$, the profile is
$(0,x^d,0,dx^{d-1})$, which gives
$\Zpoly(K_{1,d};x)=(1+x)(x^d+dx^{d-1})$.

\section{The zero-forcing polynomial of a tree}\label{sec:trees}

Let $T$ be a tree rooted at $r$.  For a vertex $v$, let $T_v$ denote the subtree induced by $v$ and its descendants, rooted at $v$.

The one-vertex rooted graph has profile
\begin{equation}\label{eq:singleton}
  (P_0,P_1,A_0,A_1)=(1,0,0,0).
\end{equation}
Indeed, the empty set is passive when the unique vertex is supplied as the root, but it cannot force an initially white root.

For each vertex $v$, recursively compute the profiles of its child subtrees.  Apply edge extension to each child profile, thereby viewing the child branch as rooted at $v$, and combine the resulting branch profiles by the rooted product.  At the global root, use \eqref{eq:normal-form}.

\begin{theorem}[Exact tree algorithm]\label{thm:tree-algorithm}
Let $T$ be a tree on $n$ vertices.  The preceding dynamic programme computes the complete zero-forcing polynomial $\Zpoly(T;x)$ exactly.  With schoolbook polynomial convolution, it uses $O(n^2)$ integer arithmetic operations on polynomial coefficients.

Let $M(\ell)$ denote the bit complexity of multiplying two $\ell$-bit integers.  A coarse implementation bound is then $O(n^2M(n))$ bit operations.  In particular, the schoolbook estimate $M(n)=O(n^2)$ gives $O(n^4)$.  The output itself has $O(n^2)$ bits, and no optimized bit-complexity claim is intended.
\end{theorem}

\begin{proof}
Correctness follows by induction on $|V(T_v)|$.  The base case is \eqref{eq:singleton}.  Assume the child profiles at $v$ are correct.  By \cref{thm:edge-extension}, edge extension gives the exact response of each child branch at $v$.  Since $T_v$ is the rooted wedge of these branches, \cref{thm:wedge} gives the exact profile at $v$.  The induction closes, and \cref{thm:normal-form} reconstructs $\Zpoly(T;x)$ at the global root.

For the arithmetic complexity, let the child-subtree sizes at a vertex $v$ be $s_1,\dots,s_d$.  Sequentially multiplying the dual-number polynomials associated with the branches requires
\[
  O\!\left(\sum_{i=1}^d s_i+\sum_{1\le i<j\le d}s_is_j\right)
\]
coefficient operations.  The product $s_is_j$ charges pairs of vertices lying in distinct child subtrees of $v$.  Each unordered pair of vertices is charged at most once, namely at its lowest common ancestor.  Summed over all $v$, these terms are therefore $O(n^2)$.

The linear edge-extension work is bounded by
\[
  \sum_{v\ne r}|V(T_v)|=O(n^2),
\]
with the path giving the extremal order.  Hence the total number of coefficient operations is $O(n^2)$.

Every coefficient is at most $2^n$, and so has $O(n)$ bits.  There are $O(n^2)$ scalar arithmetic operations on $O(n)$-bit integers, which gives the stated bit bound.  Finally, $\Zpoly(T;x)$ has $n+1$ coefficients of $O(n)$ bits each, so its binary output length is $O(n^2)$.
\end{proof}

\begin{corollary}\label{cor:forest}
The zero-forcing polynomial of an $n$-vertex forest is computable in $O(n^2)$ integer arithmetic operations on polynomial coefficients.
\end{corollary}

\begin{proof}
A set is zero forcing in a disjoint union if and only if its intersection with every connected component is zero forcing.  The polynomial of a forest is the product of the polynomials of its tree components.  Apply \cref{thm:tree-algorithm} and multiply the component polynomials.
\end{proof}

\begin{remark}
The algorithm returns the full distribution of zero-forcing sets.  In particular, it returns the zero-forcing number $Z(T)$, the number $z(T;Z(T))$ of minimum zero-forcing sets, and every higher coefficient in the same computation.
\end{remark}

\section{Transfer through blocks of bounded order}\label{sec:blocks}

We now pass from one-point unions to the block-cut tree.  A \emph{block} is a maximal connected subgraph having no cut vertex.  Every bridge is treated as a two-vertex block.  An isolated vertex is treated as a one-vertex block only when it is a connected component by itself.  If a connected graph has a single block, that block is used as the root of the block-cut tree.  Put
\[
  \beta(G)=\max\{|V(B)|:B\text{ is a block of }G\}.
\]
The purpose of this section is to prove an exact algorithm whose exponential dependence is confined to $\beta(G)$.

\subsection{Source and demand markers}

Let $(H,v)$ be a rooted graph and fix $S\subseteq V(H)\setminus\{v\}$.  Its response state is one of
\[
  P_0,\quad P_1,\quad A_0,\quad A_1.
\]
The second index records whether $S$ sources an initially white root, while the letter records whether a supplied root has no remaining obligation or must expend its unique force inside $H$.

For $\omega\in\{P_0,P_1,A_0,A_1\}$ define
\[
  \operatorname{src}(\omega)=
  \begin{cases}
  1,&\omega\in\{P_1,A_1\},\\
  0,&\omega\in\{P_0,A_0\},
  \end{cases}
\]
and
\[
  \operatorname{dem}(\omega)=
  \begin{cases}
  1,&\omega\in\{A_0,A_1\},\\
  0,&\omega\in\{P_0,P_1\}.
  \end{cases}
\]
A source marker colours the attachment vertex without external assistance.  A demand marker is a new white pendant vertex which must be forced by the attachment vertex.

\begin{lemma}[Local source--demand replacement]\label{lem:replacement}
Let $(H,v)$ be rooted, let $S\subseteq V(H)\setminus\{v\}$ have response state
$\omega$, and let $J$ satisfy $V(H)\cap V(J)=\{v\}$.  Put
$\Gamma=J\cup H$.  Form the marker graph $\Gamma_\omega$ from $J$ by adjoining
a white pendant vertex $d_v$ at $v$ when $\operatorname{dem}(\omega)=1$.
For an external initial set $Y\subseteq V(J)$, put
\[
  Y_\omega=Y\cup
  \bigl(\{v\}\text{ if }\operatorname{src}(\omega)=1\bigr).
\]
Then, for every forbidden set $Q\subseteq V(J)\setminus\{v\}$,
\begin{equation}\label{eq:replacement-projection}
  \cl_{\Gamma}^{\neg Q}(Y\cup S)\cap V(J)
  =
  \cl_{\Gamma_\omega}^{\neg Q}(Y_\omega)\cap V(J).
\end{equation}
Moreover,
\begin{equation}\label{eq:replacement-completion}
  \cl_{\Gamma}^{\neg Q}(Y\cup S)=V(\Gamma)
  \quad\Longleftrightarrow\quad
  \cl_{\Gamma_\omega}^{\neg Q}(Y_\omega)=V(\Gamma_\omega).
\end{equation}
Consequently replacement preserves the four response states and failure at
any distinguished external root $p\in V(J)\setminus\{v\}$.
\end{lemma}

\begin{proof}
Fix restricted sequences realizing $C^-(S)$ and $C^+(S)$.  If $\omega$ is
active, let $u$ be the unique vertex of
$N_H(v)\setminus C^+(S)$ and fix a restricted sequence which completes $H$
after the force $v\to u$.

We prove the two inclusions in \eqref{eq:replacement-projection} by simulating
arbitrary finite $Q$-restricted sequences.  Internal saturation steps are
allowed to be inserted during the simulation; they change no vertex of
$J\setminus\{v\}$.

\smallskip
\noindent\emph{Marker-to-branch simulation.}
Begin in $\Gamma$ by putting the $H$-part in its canonical internal
configuration.  If $v\in Y$, saturate directly from $S\cup\{v\}$ to
$C^+(S)$.  If $v\notin Y$ and $\operatorname{src}(\omega)=1$, saturate from
$S$ to $C^-(S)=C^+(S)$; this supplies $v$.  If $v\notin Y$ and the source bit
is zero, saturate from $S$ to $C^-(S)$ and leave $v$ white.

Now scan a finite legal sequence in $\Gamma_\omega$.  A force in $J$ whose
forcing vertex is different from $v$ is copied whenever its target is still
white, and is otherwise omitted.  Such a vertex has exactly the same
neighbourhood in $\Gamma$ and $\Gamma_\omega$: no vertex of
$J\setminus\{v\}$ is adjacent to $H\setminus\{v\}$ or to the demand
leaf.  Hence the copied force remains legal unless its target has already
become blue.  If a non-source branch first receives $v$ from $J$,
copy that force and then saturate \emph{directly} from $S\cup\{v\}$ to
$C^+(S)$, scanning the fixed saturation sequence from the current, possibly
larger, blue set by \cref{lem:sequence-simulation}.  This corrects the branch
state without changing $J\setminus\{v\}$.

Suppose next that the marker sequence uses $v$ as forcing vertex.  If it
forces a vertex of $J$, then no pending demand is present: while a demand leaf
is white it is an additional white neighbour of $v$, and after $v$ forces that
leaf all neighbours of $v$ are blue.  Hence the corresponding branch is
passive and has no white neighbour in $H$, so the same external force is legal
in $\Gamma$ unless its target is already blue.  If the marker force is $v\to d_v$, then all neighbours of $v$
in $J$ are blue.  In the active branch, $u$ is the unique white neighbour of
$v$ in $H$; replace $v\to d_v$ by $v\to u$ followed by the fixed restricted
completion.  Every vertex of $J$ reached by the marker sequence is
reached in $\Gamma$, and a completed marker graph expands to a completed
original graph.

\smallskip
\noindent\emph{Branch-to-marker simulation.}
Start with an arbitrary finite $Q$-restricted sequence in $\Gamma$.  First
normalize its internal chronology.  If $v\in Y$, prepend a saturation from
$S\cup\{v\}$ to $C^+(S)$.  If $v\notin Y$ and the source bit is one, prepend
the fixed saturation from $S$ to $C^-(S)=C^+(S)$, \emph{regardless of which
side coloured $v$ first in the original sequence}.  The source marker is
therefore blue from the beginning.  Scanning the original sequence from this
larger blue set and omitting forces whose targets have already become blue is
valid by \cref{lem:sequence-simulation}; in particular, a later external
force to $v$ may simply disappear.  If the source bit is zero, prepend a
saturation from $S$ to $C^-(S)$.  When an external force first colours $v$,
insert a saturation directly from $S\cup\{v\}$ to $C^+(S)$ before continuing.

After this normalization, delete every force whose two endpoints lie in
$H$, except for a force by $v$ in an active branch.  External forces with
forcing vertex different from $v$ remain legal in the marker graph for the
same neighbourhood reason: their external neighbourhoods are identical, and
a force is omitted only when its target is already blue.  An external force
by $v$ can occur only for a passive
branch: before an active branch is completed, its vertex $u$ is an additional
white neighbour of $v$, and after $v\to u$ the vertex $v$ has already used its
unique force.  No pending marker demand obstructs a copied external
force.  Finally, contract the active force $v\to u$ to the marker force
$v\to d_v$ and delete any remaining internal steps.  In the total-completion
case, the fixed internal completion may instead be appended after the
corresponding marker force.  The marker force is legal because legality of
$v\to u$ in $\Gamma$ implies that every neighbour of $v$ in $J$ is blue.
This produces a legal marker sequence reaching every vertex of $J$ reached by
the original sequence; completion of $\Gamma$ contracts to completion of
$\Gamma_\omega$.

Apply the first simulation to a maximal marker sequence and the second to a
maximal original sequence.  Confluence from \cref{lem:closure} gives the two
opposite inclusions in \eqref{eq:replacement-projection}, and the completion
statements in the two simulations give \eqref{eq:replacement-completion}.

To compare response states at an external root $p$, first take $Q=\{p\}$ and
$p\notin Y$.  Equation \eqref{eq:replacement-projection} shows that $p$ is
reached in one $p$-restricted process exactly when it is reached in the
other, so the source indicators agree even when the process later stalls.
Next take $p\in Y$ and compare both $Q=\{p\}$ and $Q=\varnothing$.
By \eqref{eq:replacement-completion} the $p$-restricted completion and the
unrestricted completion agree in the two models.  \Cref{thm:normal-form}
therefore gives the same passive, active, or
failure outcome, completing the proof.
\end{proof}

\subsection{The transfer table of a block}

Let $B$ be a block, let $r\in V(B)$ be its parent attachment, and let
\[
  U\subseteq V(B)\setminus\{r\}
\]
be the child articulation vertices.  For each $u\in U$, suppose that a descendant rooted graph has response profile
\[
  \mathbf R_u=(P_{u,0},P_{u,1},A_{u,0},A_{u,1}).
\]

For a map
\[
  \omega:U\longrightarrow\{P_0,P_1,A_0,A_1\},
\]
let $B^\omega$ be obtained from $B$ by adjoining one pendant demand marker at every $u$ for which $\operatorname{dem}(\omega(u))=1$.  For $X\subseteq V(B)\setminus\{r\}$, take
\[
  X^\omega=X\cup
  \{u\in U:\operatorname{src}(\omega(u))=1\}
\]
as the initial blue set in $B^\omega$, with $r$ excluded.  Let
\[
  \chi_{B,r}(X,\omega)
\]
denote the response category of this concrete rooted configuration.

For $\gamma\in\{P_0,P_1,A_0,A_1\}$ define
\begin{equation}\label{eq:block-transfer}
  T_\gamma(B,r;\mathbf R)
  =
  \sum_{\substack{X\subseteq V(B)\setminus\{r\}\
                   \omega\in\{P_0,P_1,A_0,A_1\}^{U}\\
                   \chi_{B,r}(X,\omega)=\gamma}}
  x^{|X|}
  \prod_{u\in U} R_{u,\omega(u)},
\end{equation}
where, for example, $R_{u,A_1}=A_{u,1}$.

\begin{theorem}[Exact block transfer]\label{thm:block-transfer}
The rooted response profile of the union of $B$ and all descendant graphs attached at the vertices of $U$, rooted at $r$, is
\[
  \boxed{
  \bigl(
    T_{P_0}(B,r;\mathbf R),
    T_{P_1}(B,r;\mathbf R),
    T_{A_0}(B,r;\mathbf R),
    T_{A_1}(B,r;\mathbf R)
  \bigr).
  }
\]
\end{theorem}

\begin{proof}
Choose an initial subset $X$ of the non-root vertices of $B$ and, independently, an initial subset in every descendant graph.  First discard descendant subsets that are blocked even when their attachment root is supplied.  Such a subset cannot occur in a globally successful configuration: if the whole configuration completed, supply that attachment vertex at time zero and delete the external step which first colours it; by monotone sequence simulation the descendant would then complete, a contradiction.  Every remaining descendant subset lies in exactly one of the four successful rooted response states.  By \cref{lem:replacement}, each such descendant graph may be replaced by its source and demand markers without changing the response category at $r$.  The resulting category is therefore $\chi_{B,r}(X,\omega)$, and the generating weight of all choices yielding the fixed pair $(X,\omega)$ is
\[
  x^{|X|}\prod_{u\in U}R_{u,\omega(u)}.
\]
The choices are disjoint and exhaustive.  Notice that a source state at $u$ and the independent choice $u\in X$ may lead to the same effective blue marker, but they are distinct original initial sets and carry weights differing by the factor $x$; the explicit factor $x^{|X|}$ records this distinction.  Summing according to the response category gives \eqref{eq:block-transfer}.
\end{proof}

The theorem gives a finite kernel for every block.  It is deliberately stated for an arbitrary block rather than only for a clique or a cycle.

\subsection{Effective four-letter aggregation}\label{subsec:effective-block}

The transfer table above keeps the choice of a block vertex separate from the
four-state response of its descendant.  For computation these choices can be
aggregated without losing information.  At a non-root child articulation $u$
define
\begin{align}
 \lambda_u^{00}&=P_{u,0},
 &\lambda_u^{10}&=xP_{u,0}+(1+x)P_{u,1},
 \label{eq:block-letter-a}\\
 \lambda_u^{01}&=A_{u,0},
 &\lambda_u^{11}&=xA_{u,0}+(1+x)A_{u,1}.
 \label{eq:block-letter-b}
\end{align}
The first bit is the effective initial colour seen by the block and the second
is the pending-demand bit.  For example, a source state $P_1$ produces
an effectively blue articulation whether or not the articulation itself was
selected; the two disjoint choices have total weight $(1+x)P_{u,1}$.  At a
non-articulation block vertex the only possible letters are $00$ and $10$,
with weights $1$ and $x$.

\begin{theorem}[Effective arbitrary-block kernel]\label{thm:effective-block}
Let $B$ have order $m$, parent root $r$, and $q$ child articulations.  For each
assignment of one effective letter to every vertex of $V(B)\setminus\{r\}$,
adjoin a demand leaf at every letter with second bit one and initially colour
every vertex whose first bit is one.  Classify the resulting concrete rooted
marker graph as $P_0,P_1,A_0,A_1$, or failure, and sum the product of its local
letter weights over each successful class.  The four sums are exactly the
response profile of the decorated block.

The number of concrete marker configurations is at most
\[
  4^q2^{m-1-q}=2^{m-1+q}\le4^{m-1}.
\]
\end{theorem}

\begin{proof}
Expand each local weight in \eqref{eq:block-letter-a}--\eqref{eq:block-letter-b}.
Its summands are precisely the disjoint choices of descendant response state
and selection/nonselection of the articulation that produce the same effective
blue/demand pair.  By \cref{lem:replacement}, all summands aggregated into a
fixed letter have identical behaviour inside the parent block.  Distributivity
therefore turns the sum in \cref{thm:block-transfer} into the stated four-letter
enumeration, with no loss or duplication of original initial sets.  Ordinary
block vertices contribute only the two letters $00$ and $10$.
\end{proof}

\subsection{Clique blocks of arbitrary order}\label{subsec:clique-transfer}

For a clique block the effective configurations admit a closed transfer.  Let
$B=K_m$ have parent root $r$.  At each non-root vertex $v_i$, use the letter
weights \eqref{eq:block-letter-a}--\eqref{eq:block-letter-b}, taking the
singleton profile $(1,0,0,0)$ when no descendant is attached.  Put
\[
  B_i=\lambda_i^{10}+\lambda_i^{11},\qquad
  D_i=\lambda_i^{11},\qquad
  W_i=\lambda_i^{00}+\lambda_i^{01},
\]
and define
\begin{align*}
  B_*&=\prod_i B_i,
  &D_*&=\prod_i D_i,\\
  W_*&=\sum_i W_i\prod_{j\ne i}B_j,
  &E_*&=\sum_i W_i\prod_{j\ne i}D_j.
\end{align*}
The term $B_*$ counts configurations in which every non-root clique vertex is
effectively blue, $D_*$ the subfamily in which each such vertex also carries
a demand, $W_*$ configurations with exactly one effective white vertex, and
$E_*$ the latter subfamily in which every effective blue vertex carries a
demand.

\begin{theorem}[Clique transfer]\label{thm:clique-transfer}
For a clique block of arbitrary order $m\ge2$, the response profile at its
parent root is
\[
  \boxed{
  P_0=D_*+W_*-E_*,\qquad
  P_1=B_*-D_*,\qquad
  A_0=E_*,\qquad
  A_1=0.
  }
\]
The four aggregates can be computed without division in one scan.  Starting
with $B_*=D_*=1$ and $W_*=E_*=0$, at vertex $i$ update simultaneously
\begin{align*}
  B_*&\leftarrow B_*B_i,
  &D_*&\leftarrow D_*D_i,\\
  W_*&\leftarrow W_*B_i+B_*W_i,
  &E_*&\leftarrow E_*D_i+D_*W_i,
\end{align*}
using the old values on the right.
\end{theorem}

\begin{proof}
A blue clique vertex carrying a white demand leaf cannot force inside the
clique: until its demand is discharged the leaf is an additional white
neighbour, and the demand can be discharged only after every clique neighbour
is blue.  Demand vertices are prohibited from supplying the unique force
that colours a remaining white clique vertex.

With the root supplied, two or more effectively white non-root vertices block
all forces into the clique, so only the cases counted by $B_*$ and $W_*$ can
succeed.  If all non-root vertices are effectively blue, the root-forbidden
process completes.  It has source bit one exactly when at least one effective
blue vertex has no demand, yielding $P_1=B_*-D_*$ and the all-demand remainder
$D_*$ in $P_0$.

If exactly one non-root vertex is effectively white, then with an initially
white root there are two white clique vertices and no source force can begin,
so the source bit is zero.  With the root supplied, a blue non-demand vertex
can force the unique white clique vertex unless every effective blue vertex
has a demand.  The former configurations have weight $W_*-E_*$ and are
passive; in the latter, of weight $E_*$, the root must make the unique clique
force and the state is active.  Once the clique is blue all demands can be
discharged.  An active source state $A_1$ is therefore impossible.  The scan
recurrences are the standard division-free updates for the displayed products
and one-exception sums.
\end{proof}

\subsection{The block-cut dynamic programme}

Let $G$ be connected.  Root its block-cut tree at a leaf block and choose a
non-articulation vertex of that block as the global root.  Process nodes from
the leaves towards the root.  At an articulation node combine all child-block
profiles by \cref{thm:wedge}.  At a block node apply
\cref{thm:block-transfer}.  At the global root apply
\cref{thm:normal-form}.  For a disconnected graph, multiply the polynomials of
its connected components.

We use the following accounting convention.  A profile returned across an
attachment vertex represents only the vertices strictly below that attachment;
the attachment vertex itself is excluded and is assigned to the parent block.
At an articulation node, the represented sets of the child blocks are
pairwise disjoint.  At a block with parent attachment $r$, the represented set
is the disjoint union of $V(B)\setminus\{r\}$ and the represented sets returned
by its child articulations.  This convention also explains why each initial
selection factor $x$ is introduced exactly once.

\begin{lemma}[Pair charging on a decomposition tree]\label{lem:pair-charging}
Let a rooted decomposition tree have $O(n)$ nodes, with an absolute constant.
At every node, partition the newly combined represented set into pairwise
disjoint pieces of sizes $s_1,\ldots,s_q$, and suppose the node returns its
entire represented set as one piece to its parent.  Suppose that the polynomial
work at that node is
\[
  O\!\left(\sum_i s_i+\sum_{i<j}s_is_j\right).
\]
If pieces are merged but never split on the path to the root, then the total
work is $O(n^2)$ coefficient operations.
\end{lemma}

\begin{proof}
Charge each occurrence in a linear term to the ordered pair consisting of the
represented vertex and the current decomposition node.  Because the
decomposition has $O(n)$ nodes, there are $O(n^2)$ such pairs.  Charge a term $s_is_j$ to unordered pairs of original
vertices whose endpoints lie in the two corresponding pieces.  Each unordered
pair is charged at the unique lowest node at which its endpoints lie in
distinct pieces.  At every strict ancestor the two endpoints already belong to
the same merged piece, so the pair is not charged again.  There are at most
$\binom n2$ unordered pairs.
\end{proof}

\begin{theorem}[Bounded-block algorithm]\label{thm:bounded-block}
Let $G$ have $n$ vertices and suppose every block has at most $b$ vertices.
The transfer tables can be constructed using
\[
  O(4^{b}b^2n)
\]
elementary combinatorial operations.  After this preprocessing,
$\Zpoly(G;x)$ is computable using
\[
  \boxed{O(4^{b}n^2)}
\]
integer arithmetic operations on polynomial coefficients under schoolbook
polynomial convolution.  In particular, for fixed $b$ the complete polynomial is
computable in polynomial time.
\end{theorem}

\begin{proof}
Correctness follows by induction on the rooted block-cut tree from
\cref{thm:wedge,thm:block-transfer,thm:normal-form}.

For a block $B$ with parent attachment $r$ and $q$ child articulations, there
are at most
\[
  4^q2^{|V(B)|-1-q}\le 4^{b-1}
\]
effective marker configurations by \cref{thm:effective-block}.  Each response category is determined by a constant number of restricted-closure
computations on a graph with at most $2b-1$ vertices.  With an adjacency
matrix, or adjacency lists together with a queue of currently legal forcing
vertices, each such computation uses $O(b^2)$ elementary graph operations.
This gives the stated combinatorial preprocessing bound over $O(n)$ blocks.

For a fixed table entry, multiplying one chosen polynomial from child profiles
of represented sizes $s_1,\ldots,s_q$ has the cost in
\cref{lem:pair-charging}.  Multiplying by the number of table entries gives the
factor $O(4^{b})$; the wedge operations at articulation nodes have the same
form.  The lemma yields $O(4^{b}n^2)$ coefficient operations.
\end{proof}

\begin{remark}[Arithmetic model and bit complexity]\label{rem:bit-complexity}
Every coefficient is at most $2^n$ and therefore has $O(n)$ bits.  The theorem
is a unit-cost count of integer arithmetic operations on polynomial
coefficients, not a RAM-time bound.
With schoolbook integer multiplication, the preceding implementation has the
coarse bit bound $O(4^{b}n^4)$, in addition to the combinatorial preprocessing.
No optimal bit-complexity claim is intended.
\end{remark}

\section{An explicit cycle automaton}\label{sec:cactus}

Finite automata have previously been associated with zero forcing: the
zero-forcing automata of Shamsizadeh, Zahedi, Golmohamadian, and Abolpour, and
the later $k$-forcing automata, use the vertices of the input graph as automaton
states and encode a chosen forcing process and its recognized language
\cite{ShamsizadehEtAl2021,ShamsizadehEtAl2024}.  The construction below is of
a different kind.  It is a fixed-size fort-frontier recognizer whose input is
a color--demand word around a cycle; its graph-independent five frontier
states generate a $100$-element Boolean transition monoid and a certified
$32$-state deterministic response quotient.  The claim here is the
specific finite frontier representation and its transfer consequences, not
the introduction of automata into zero forcing.

A \emph{cactus} is a graph in which any two cycles have at most one vertex in common; equivalently, every edge belongs to at most one cycle.  Every block of a nontrivial connected cactus is either a bridge or a cycle; isolated vertices are handled separately.  The bounded-block algorithm is exponential in the order of a block, and the lengths of the
cycle blocks need not be bounded.  We now exploit their cyclic order.  The result is
a constant-state transfer algorithm which is independent of the cycle length.

\subsection{Restricted forts}

Let $H$ be a graph and let $Q\subseteq V(H)$.  A forcing process is
\emph{$Q$-restricted} if no vertex of $Q$ is used as a forcing vertex.

\begin{definition}
A non-empty set $F\subseteq V(H)$ is a \emph{$Q$-fort} if
\[
  |N_H(v)\cap F|\ne 1
  \qquad
  \text{for every }v\in V(H)\setminus(F\cup Q).
\]
The usual forts are the $\varnothing$-forts.
\end{definition}

\begin{lemma}[Restricted fort criterion]\label{lem:q-fort}
Let $B\subseteq V(H)$ and let $t\in V(H)$.
\begin{enumerate}[label=\textnormal{(\roman*)}]
\item The $Q$-restricted process from $B$ colours all of $H$ if and only if
      $B$ intersects every $Q$-fort.
\item The vertex $t$ belongs to the $Q$-restricted closure of $B$ if and only
      if $B$ intersects every $Q$-fort containing $t$.
\end{enumerate}
\end{lemma}

\begin{proof}
Suppose that a $Q$-fort $F$ is initially white.  If a first force entered
$F$, its forcing vertex would lie outside $F\cup Q$ and would have exactly one
white neighbour in $F$, contrary to the definition.  Hence no vertex of $F$
is ever coloured.

Conversely, let $W$ be the set of vertices left white at the terminal
restricted closure.  Every blue vertex outside $Q$ has either no neighbour in
$W$ or at least two, since otherwise one further force would be available.
Whenever $W\ne\varnothing$, it is a $Q$-fort disjoint from $B$.  If a
specified vertex $t$ remains white, then $t\in W$.  These observations prove
both statements.
\end{proof}

Consider a cycle block
\[
  C=v_0v_1\cdots v_{m-1}v_0,
  \qquad v_0=r,
\]
where $r$ is the parent attachment.  After applying the replacement theorem
of \cref{lem:replacement}, a non-root cycle vertex may carry a pending demand
leaf.

\begin{lemma}[Demand suppression]\label{lem:demand-suppression}
A cycle vertex carrying a pending demand leaf never forces a cycle neighbour.
Moreover, once all cycle vertices are blue, every pending demand can be
completed.  Consequently, for the purpose of deciding whether the marked
cycle completes, a demand vertex may be treated as a vertex forbidden from
forcing along the cycle.
\end{lemma}

\begin{proof}
Suppose that a demand vertex $v$ forces a cycle neighbour at some first such
step.  Its demand leaf must already be blue; otherwise the leaf and the target
cycle neighbour would be two white neighbours of $v$.  The leaf can only have
been forced by $v$, and that earlier force required both cycle neighbours of
$v$ to be blue.  This contradicts the choice of the later white target.

Conversely, if all cycle vertices are blue, then each demand vertex has its
leaf as its unique white neighbour and may force it.  The demands are
independent, so all of them can be completed.
\end{proof}

\subsection{A five-state fort frontier}

For $1\le i<m$, let $b_i$ record effective initial blueness at $v_i$, and let
$d_i$ record a pending demand.  For a candidate fort $W$, put
$w_i=\one_{\{v_i\in W\}}$.  At a non-root vertex the disjointness and fort
conditions are
\begin{equation}\label{eq:local-fort-condition}
  b_iw_i=0,
  \qquad
  w_i=0\text{ and }d_i=0\Longrightarrow w_{i-1}=w_{i+1}.
\end{equation}

Before processing the letter at $v_i$, a frontier state is
\[
  (w_{i-1},w_i,h_i),
  \qquad
  h_i=\bigvee_{j=0}^{i}w_j.
\]
After choosing $w_{i+1}$ and checking \eqref{eq:local-fort-condition} at
$v_i$, the state becomes
\[
  (w_i,w_{i+1},h_i\vee w_{i+1}).
\]
The only reachable triples are
\[
 q_0=(0,0,0),\ q_1=(0,0,1),\ q_2=(0,1,1),\
 q_3=(1,0,1),\ q_4=(1,1,1).
\]
There is no sixth state because $h_i=0$ forces both displayed fort bits to be
zero.

For a letter $(b,d)$ define the Boolean $5\times5$ matrix $T_{bd}$ by
\begin{equation}\label{eq:letter-transition}
  T_{bd}\bigl((a,c,h),(c,z,h\vee z)\bigr)=1
\end{equation}
exactly when $bc=0$ and
$c=0,d=0\Longrightarrow a=z$.  Row vectors act on the left, so the word
$(b_1,d_1)\cdots(b_{m-1},d_{m-1})$ is represented by
\[
  M=T_{b_1d_1}\cdots T_{b_{m-1}d_{m-1}}.
\]
An entry $M_{ab}=1$ records a trace from initial frontier state $q_a$ to final
state $q_b$.  The initial row encodes $(w_0,w_1,w_0\vee w_1)$; after the last
letter the final column encodes $(w_{m-1},w_0,h)$, because closing the cycle
sets the next bit equal to $w_0$.

Let $E_R,E_O,E_S$ denote, respectively, an obstruction with the root blue and
forbidden, an obstruction with the blue root allowed to force, and an
obstruction containing an initially white root.  In the order
$q_0,\ldots,q_4$,
\begin{align}
 E_R(M)&\Longleftrightarrow
     \bigvee_{a\in\{0,2\}}\ \bigvee_{b\in\{1,3\}}M_{ab}=1,
     \label{eq:ER}\\
 E_O(M)&\Longleftrightarrow M_{01}\vee M_{23}=1,
     \label{eq:EO}\\
 E_S(M)&\Longleftrightarrow
     \bigvee_{a\in\{3,4\}}\ \bigvee_{b\in\{2,4\}}M_{ab}=1.
     \label{eq:ES}
\end{align}
For $E_R$ the root bit is $0$ and the final history bit enforces non-emptiness.
For $E_O$ the root lies outside the fort and is allowed to force, so its two
cycle neighbours must have equal fort bits.  Every ordinary fort counted by
$E_O$ is also a root-restricted fort counted by $E_R$; hence
$E_O\Rightarrow E_R$, and the response cases below are disjoint and
exhaustive.  For $E_S$ the root bit is $1$; non-emptiness is automatic and the
fort condition at the root is omitted because the source test is carried out
in the root-forbidden process, even when the root itself belongs to the fort.

\begin{example}[An unmarked triangle]\label{ex:triangle-automaton}
For a triangle with no selected non-root vertex and no demand, the word is
$00\,00$.  The product $T_{00}^2$ has, among others, the transition
$q_0\to q_0$, corresponding to the empty trace, and transitions
$q_2\to q_2,q_3,q_4$, corresponding to non-empty fort traces.  In particular
$E_O(T_{00}^2)$ holds, so the readout is failure: a supplied root alone does
not zero force a triangle.  If both non-root vertices are selected, the word
is $10\,10$; the readout is $P_1$, as the selected vertices force an initially
white root and the cycle is already complete when the root is supplied.
\end{example}

\begin{proposition}[Response readout]\label{prop:cycle-readout}
The concrete response of the marked cycle word $\mathbf a$ is determined by
$M=M(\mathbf a)$ as follows:
\[
\begin{array}{c|c}
\text{condition}&\text{response}\\
\hline
E_O(M)&\text{failure},\\
\neg E_R(M),\ \neg E_S(M)&P_1,\\
\neg E_R(M),\ E_S(M)&P_0,\\
E_R(M),\ \neg E_O(M),\ \neg E_S(M)&A_1,\\
E_R(M),\ \neg E_O(M),\ E_S(M)&A_0.
\end{array}
\]
\end{proposition}

\begin{proof}
By \cref{lem:demand-suppression,lem:q-fort}, the plus process with the root
forbidden completes precisely when $E_R(M)$ is false.  This is the passive
case.  With the root allowed, the plus process completes precisely when
$E_O(M)$ is false.  By Theorem~\ref{thm:normal-form}, the response is active exactly when $E_R(M)$ is true and $E_O(M)$ is false: the root-forbidden saturation fails, but one root force followed by another restricted saturation succeeds.  If $E_O(M)$ is true, even the unrestricted plus process
fails.  Finally, the source bit is one precisely when there is no restricted
fort containing the root, that is, when $E_S(M)$ is false.
\end{proof}

The four matrices $T_{00},T_{10},T_{01},T_{11}$ generate a finite Boolean
transition monoid.  The certificate described in
\cref{app:automaton-certificate} lists all $100$ matrices in canonical
breadth-first order and records their four right transitions and response
outputs.  Response equivalence identifies matrices having the same output
after every continuation word.  The supplied quotient map has $32$ classes.
For every pair of distinct quotient states, the certificate supplies a
continuation word on which their outputs differ; hence the quotient is minimal
as a deterministic response automaton.  Minimality is not used for correctness
or complexity--the unreduced $100$-state monoid already suffices--but the
smaller quotient is used in the reference implementation.

\subsection{Weighted cycle transfer}

Let $\mathbf R_i=(P_{i,0},P_{i,1},A_{i,0},A_{i,1})$ be the descendant response
profile at $v_i$; at a non-articulation vertex use the singleton profile
$(1,0,0,0)$.  Summing over whether $v_i$ itself is selected gives the four
letter weights
\begin{align}
 \lambda_i^{00}&=P_{i,0},
 &\lambda_i^{10}&=xP_{i,0}+(1+x)P_{i,1},
 \label{eq:cycle-weights-a}\\
 \lambda_i^{01}&=A_{i,0},
 &\lambda_i^{11}&=xA_{i,0}+(1+x)A_{i,1}.
 \label{eq:cycle-weights-b}
\end{align}
The first coordinate is the effective blue bit and the second is the demand
bit.  The aggregation is displayed explicitly below; ``selected'' contributes the
factor $x$.
\[
\begin{array}{c|c|c|c}
\text{descendant state}&v_i\text{ selected?}&(b_i,d_i)&\text{weight}\\
\hline
P_0&0&00&P_{i,0}\\
P_0&1&10&xP_{i,0}\\
P_1&0&10&P_{i,1}\\
P_1&1&10&xP_{i,1}\\
A_0&0&01&A_{i,0}\\
A_0&1&11&xA_{i,0}\\
A_1&0&11&A_{i,1}\\
A_1&1&11&xA_{i,1}
\end{array}
\]
The rows are disjoint at the level of original initial sets even when they
produce the same effective letter, so no double counting occurs.

Let $\mathcal A$ be either the $100$-state transition monoid or the $32$-state
quotient.  Starting from the identity state $e$, define polynomial weights
$f_i(s)$ by
\[
  f_0(e)=1,
  \qquad
  f_0(s)=0\quad(s\ne e),
\]
and
\begin{equation}\label{eq:cycle-dp}
  f_i(t)=
  \sum_{\substack{s\in\mathcal A,
                   (b,d)\in\{0,1\}^2\\
                   s\cdot(b,d)=t}}
  f_{i-1}(s)\lambda_i^{bd}.
\end{equation}
At the end, sum $f_{m-1}(s)$ over the states whose readout is respectively
$P_0,P_1,A_0,A_1$.

\begin{theorem}[Explicit cycle transfer]\label{thm:cycle-transfer}
The four polynomials obtained from \eqref{eq:cycle-dp} are exactly the rooted
response profile of the cycle block together with all descendant graphs
attached at its non-root vertices.
\end{theorem}

\begin{proof}
For every non-root cycle vertex, the four weights
\eqref{eq:cycle-weights-a}--\eqref{eq:cycle-weights-b} form a disjoint and
exhaustive partition of the choices of the block vertex and its descendant
response state.  Their product is the generating weight of the corresponding
marked word.  The deterministic transition in \eqref{eq:cycle-dp} computes its
fort matrix, or its response-equivalence class.  By
\cref{prop:cycle-readout}, the terminal readout is exactly its rooted response
category.  Summing over all words proves the assertion.
\end{proof}

\subsection{The cactus algorithm}

A cactus is processed on its rooted block-cut tree.  A connected component
consisting of one isolated vertex is a base case: its polynomial is $x$, and
its rooted profile is $(1,0,0,0)$.  Every other connected component is
processed as follows.  The description is also pseudocode for the algorithm.
\begin{enumerate}[label=\textnormal{(\arabic*)}]
\item In each connected component choose a leaf block, and choose a
non-articulation vertex $r$ of that block as global root.
\item Process the block-cut tree in postorder.
\item At an articulation $v$, combine all already processed child-block
profiles by the rooted product.  The resulting profile represents the entire
descendant graph attached at $v$; the initial choice of $v$ itself is not yet
counted.
\item If a block is a bridge $rv$, apply the pendant-edge transfer to move the
profile at $v$ to the parent attachment $r$.
\item If a block is a cycle, orient it from its parent attachment
$v_0=r$ as $v_0v_1\cdots v_{m-1}v_0$.  Insert at each $v_i$, $i\ge1$, the
profile obtained in step~(3), or the singleton profile when no child block is
present.  Apply the weighted cycle transfer.  The factors $x$ in
\eqref{eq:cycle-weights-a}--\eqref{eq:cycle-weights-b} count the cycle vertex
$v_i$ exactly once.
\item Return the resulting block profile to the parent articulation.  At the
global root apply \cref{thm:normal-form}; multiply component polynomials for a
disconnected cactus.
\end{enumerate}

\begin{theorem}[Quadratic cactus algorithm]\label{thm:cactus}
The complete zero-forcing polynomial of an $n$-vertex cactus graph can be
computed exactly with $O(n)$ graph-structural preprocessing and
\[
  \boxed{O(n^2)}
\]
integer arithmetic operations on polynomial coefficients under schoolbook
polynomial convolution.  The
algorithm permits cycles of arbitrary length.  A coarse schoolbook bit bound
is $O(n^4)$.
\end{theorem}

\begin{proof}
Correctness follows by postorder induction from
\cref{thm:wedge,thm:edge-extension,thm:cycle-transfer}.  Every non-root graph
vertex is counted exactly once: at a bridge transfer if it is the child
endpoint of a bridge block, or in one of the eight rows of the cycle-letter
table if it lies in a cycle block.  Articulation products combine descendant
choices but introduce no new factor $x$.

The block-cut tree and cyclic orders are found in linear time.  The automaton
has constant size.  At a cycle block, let $s_i$ be one plus the number of
vertices represented by the descendant profile inserted at the non-root cycle
vertex $v_i$.  The added one is $v_i$ itself; the parent attachment $v_0$ is
excluded.  These pieces are pairwise disjoint under the convention preceding
\cref{lem:pair-charging}.  The scan costs
\[
  O\!\left(\sum_i s_i+\sum_{i<j}s_is_j\right)
\]
coefficient operations, up to an absolute automaton constant.  Wedge products
and bridge transfers have the same form.  Apply
\cref{lem:pair-charging} to obtain $O(n^2)$ coefficient operations.  The bit
bound follows from the $O(n)$-bit coefficient bound as in
\cref{rem:bit-complexity}.
\end{proof}

\begin{theorem}[Clique--cycle block algorithm]\label{thm:clique-cycle}
Let $G$ be an $n$-vertex graph every block of which is either a clique or a
cycle (with $K_2$ serving as a bridge block).  Then the complete zero-forcing
polynomial can be computed exactly after $O(n+|E(G)|)$ structural preprocessing
using
\[
  \boxed{O(n^2)}
\]
integer arithmetic operations on polynomial coefficients under schoolbook
polynomial convolution.  The block orders and cycle lengths are unrestricted.
Consequently this gives an explicit quadratic coefficient-operation algorithm
for every block graph as well as every cactus, and more generally for every
block-cactus graph.  The class has unbounded treewidth, since it contains
$K_m$ for every $m$, but it has bounded clique-width; see
\cite{Gurski2017} and the discussion in \cref{sec:treewidth}.
\end{theorem}

\begin{proof}
Build the block-cut forest in $O(n+|E(G)|)$ time and process it in postorder.
At articulation nodes use \cref{thm:wedge}; at clique blocks use
\cref{thm:clique-transfer}; and at non-clique cycle blocks use
\cref{thm:cycle-transfer}.  (The edge $K_2$ may equivalently be handled by
\cref{thm:edge-extension}.)  Each block transfer performs only a constant
number of polynomial convolutions per represented piece.  The represented
interiors returned by distinct children are disjoint, so the charging argument
of \cref{lem:pair-charging} gives $O(n^2)$ coefficient operations globally.
For dense clique blocks the input itself may have $\Theta(n^2)$ edges, which is
why the structural preprocessing is stated as $O(n+|E(G)|)$ rather than
$O(n)$.
\end{proof}

\subsection{Independent random initial sets and vertex influence}\label{subsec:influence}

The same transfer is naturally multivariate.  Let every vertex $v$ be selected
independently with probability $p_v$, and write
\[
 Q_G(\mathbf p)=\Pr_{\mathbf p}(S\text{ is a zero-forcing set of }G).
\]
Equivalently, one may regard
$F_G(\mathbf y)=\sum_{S\text{ zero forcing}}\prod_{v\in S}y_v$ as a
multivariate generating function represented by a division-free arithmetic
circuit rather than by its exponentially large expanded coefficient table.

For a non-root block vertex $v_i$ carrying descendant profile
\[
 (P_{i,0},P_{i,1},A_{i,0},A_{i,1}),
\]
the four effective source--demand letter weights become
\begin{align}
 \lambda_i^{00}&=(1-p_i)P_{i,0}, &
 \lambda_i^{10}&=p_iP_{i,0}+P_{i,1},\label{eq:prob-letters-a}\\
 \lambda_i^{01}&=(1-p_i)A_{i,0}, &
 \lambda_i^{11}&=p_iA_{i,0}+A_{i,1}.\label{eq:prob-letters-b}
\end{align}
At the final root $r$ the readout is
\begin{equation}\label{eq:prob-readout}
 Q_G(\mathbf p)=p_r(P_0+A_0)+P_1+A_1.
\end{equation}

\begin{corollary}[All vertex influences in linear scalar arithmetic]
\label{cor:influences}
On a cactus, or more generally on a graph whose blocks are cliques or cycles,
with its block structure supplied, $Q_G(\mathbf p)$ and all first partial
derivatives $\partial Q_G/\partial p_v$ can be evaluated using $O(n)$ scalar
arithmetic operations.  If the structure is not supplied, add
$O(n+|E(G)|)$ preprocessing.
Moreover,
\[
 \frac{\partial Q_G}{\partial p_v}
 =\Pr_{\mathbf p_{-v}}(S\cup\{v\}\text{ succeeds})
  -\Pr_{\mathbf p_{-v}}(S\text{ succeeds}),
\]
so by monotonicity this derivative is the probability that $v$ is pivotal.
\end{corollary}

\begin{proof}
Replace every factor $x$ and every unselected weight $1$ in the clique/cycle
transfer by the corresponding probabilities as in
\eqref{eq:prob-letters-a}--\eqref{eq:prob-letters-b}.  The postorder algorithm
then forms a division-free straight-line circuit with $O(n)$ addition and
multiplication gates and terminal readout \eqref{eq:prob-readout}.  The
Baur--Strassen theorem gives all first partial derivatives of an arithmetic
circuit with only constant-factor overhead \cite{BaurStrassen1983}; equivalently,
one may run the usual reverse-mode adjoint sweep on this particular circuit.
All derivatives require $O(n)$ scalar operations.  The displayed identity
follows by conditioning on whether $v$ is selected; multilinearity makes the
derivative the difference of the two conditional success probabilities.
Zero-forcing success is increasing under enlarging the initial set, so the
difference is the pivotal probability.  The contribution specific to the
present setting is the explicit division-free $O(n)$-gate heterogeneous
success circuit, not the general reverse-mode differentiation principle.
\end{proof}

\begin{corollary}[Zero influence and fort-irrelevant vertices]
\label{cor:zero-influence-fort}
For any finite graph $G$, fix $v\in V(G)$ and suppose $0<p_u<1$ for every
$u\ne v$.  Then
\[
\begin{aligned}
 \frac{\partial Q_G}{\partial p_v}=0
 &\quad\Longleftrightarrow\quad
 v\text{ belongs to no inclusion-minimal zero-forcing set},\\
 &\quad\Longleftrightarrow\quad
 v\text{ belongs to no inclusion-minimal fort}.
\end{aligned}
\]
\end{corollary}

\begin{proof}
Conditioning on the membership of $v$ in the product measure, or simply
differentiating the multilinear polynomial $Q_G$, shows that the derivative is
the probability, over subsets of $V(G)\setminus\{v\}$, that adding $v$
changes failure into success.  Full
support on the other vertices makes this probability zero exactly when no
subset $T\subseteq V(G)\setminus\{v\}$ is failed while $T\cup\{v\}$ is
zero forcing.  Such a witness exists exactly when $v$ belongs to an
inclusion-minimal zero-forcing set: minimize $T\cup\{v\}$ among its
zero-forcing subsets containing $v$; conversely removing $v$ from a minimal
zero-forcing set must fail.  Jacob's irrelevance theorem identifies the
vertices lying in no minimal zero-forcing set with those lying in no minimal
fort \cite[Theorem~2.4]{Jacob2025}.  The last equivalence is an application
of that established characterization, not a new characterization of
irrelevant vertices.
\end{proof}

\section{Logical bounded-width context}\label{sec:treewidth}

The zero-forcing predicate is $\mathrm{MSO}_1$-definable: with vertex-set
variables $S$ and $F$, one says that $F$ is a nonempty fort and that $S$ meets
every fort.  Standard counting metatheorems therefore evaluate
$\Zpoly(G;x)$ on effectively bounded clique-width classes when an appropriate
expression is supplied; fixed-treewidth consequences follow from the usual
MSO framework \cite{ArnborgLagergrenSeese1991,CourcelleMakowskyRotics2001}.
These observations are background, not tractability claims of this paper.

They are relevant to the block-cactus application.  Gurski records that the
clique-width of a graph is at most two more than the maximum clique-width of a
block \cite[Sec.~3.6]{Gurski2017}.  Since cliques and cycles have bounded
clique-width, the class treated in \cref{thm:clique-cycle} is already covered
by the general metatheory.  What \cref{thm:clique-cycle} adds is an explicit
small-state construction, a direct implementation, and an audited
$O(n^2)$ coefficient-operation bound.  Recent parameterized work on minimum
zero-forcing sets and related variants concerns a different optimization or
decision problem \cite{BhyravarapuEtAl2025,GoettlicherEtAl2026}.

\section{Computational verification}\label{sec:verification}

The computations serve as checks or as finite certificates; none is used in
place of an infinite proof.  Independent subset enumeration was used wherever
practical.  The main checks are summarized below; the supplement contains the
full logs, seeds, expected outputs, and commands.

\begin{center}
\begin{tabular}{@{}p{0.64\textwidth}p{0.28\textwidth}@{}}
\toprule
Check & Cases or certificate size \\
\midrule
Root normal form, edge transfer, wedges, articulation cases, and trees
  & $34{,}625$ \\
Atlas and random bounded-block, cactus, and block-cactus comparisons
  & $2{,}583$ \\
Marked cycle and clique words
  & $174{,}756$ \\
Projected-closure replacement and decorated clique/block checks
  & $69{,}966$ \\
Separator compatibility checks and realized chain states
  & $8{,}284$ \\
Cycle certificate
  & $100$ matrices, $32$ states, $496$ witnesses \\
Two-terminal separator-rank certificate
  & ranks $52,44,16,60$ \\
\bottomrule
\end{tabular}
\end{center}

The separator verifier is independent of the block and cycle code.  It builds
the atom posets and chain states from the definitions, checks all $52$
two-terminal chain realizations by direct prefort enumeration, samples the
three-terminal construction, and reconstructs the $q=2$ coefficient matrices.
Exact rational elimination gives the symbolic upper ranks.  Specialization at
$x=2$, followed by elimination modulo $1{,}000{,}003$, gives matching lower
ranks.

The cycle verifier reconstructs the four $5\times5$ frontier matrices,
regenerates the $100$-element Boolean transition monoid, checks the quotient to
$32$ response states, and verifies a distinguishing continuation for each of
the $496$ unordered pairs.  The full audit also reconstructs the rank-four
connection minor by direct subset enumeration and obtains
$4x^6(x+1)^2$.

\section{Discussion}\label{sec:discussion}

The separator theorem gives an explicit interface for every fixed boundary
size.  Its three ingredients are the prefort atoms, the compatibility rule at
the boundary, and the chain compression.  The realization theorem is what
turns the last step from an upper-bound construction into an exact connection
rank statement.  The four-state rooted calculus is the case $q=1$.

For independent boundaries the dimension is explicit:
$r_q^{\circ}=c_q$.  Boundary edges add Boolean M\"obius sectors, leading to
\eqref{eq:master-rank}; already $c_2=52$ and $r_2=60$.  A natural next problem
is to understand the ranks of those sectors in closed form or asymptotically.
It would also be useful to find presentations of the higher connection
algebras comparable to the source/capacity presentation at $q=1$, and to see
which graph families force similarly small quotients of the universal
interface.

The explicit one-terminal algorithms have a different purpose.  The general
separator algebra need not be an efficient implementation when $q$ is large,
whereas trees, cliques, and cycles collapse to very small concrete states.
This is why the structural separator theory and the block algorithms are kept
in the same paper.

We use existing machinery where it is appropriate.  Connection matrices and
graph algebras are classical \cite{Lovasz2006,KotekMakowsky2014}; forts and
fort transversals are standard in zero forcing
\cite{AIM2008,CameronHogbenEtAl2026}; minimal-fort enumeration is treated in
\cite{BeckerEtAl2025,CameronLi2025,DatKenter2026,AbiadGhasemiNezhad2026}; and
finite automata for forcing processes occur in
\cite{ShamsizadehEtAl2021,ShamsizadehEtAl2024}.  The claims made here concern
the explicit separator interface, its chain quotient and realization, the
resulting exact connection-rank formula, and the full-polynomial transfer
systems.  We do not claim priority for the general existence of connection
matrices, forts, automata, bounded-width tractability, or automatic
differentiation.

\appendix
\section{Finite certificates}\label{app:automaton-certificate}

Two finite claims in the paper are supplied with independent certificates.
The file \texttt{cycle\_automaton\_certificate.json} lists the generated
$5\times5$ Boolean matrices, the quotient map to the $32$ response states,
and a distinguishing continuation for each unordered pair of quotient states.
The verifier rebuilds the monoid from the four generators in
\eqref{eq:letter-transition} before checking the listed data.

The file \texttt{separator\_rank\_certificate.json} records the exact
coefficient-span and specialization data used for the $q=2$ calculation.  Its
verifier reconstructs the finite kernels, checks the stated minors, and
repeats the modular rank calculation.  Both verifiers are independent of the
main dynamic-programming implementations.  The checksum manifest in the
supplement identifies the files used for the reported calculations.

\section*{Data and code availability}

The supplementary archive contains the source code, the two certificate files
and their verifiers, the separator-realization audits, the rooted/tree/block/
cycle implementations, pinned dependencies, expected outputs, and a checksum
manifest.  No external data set is used.  The complete audit is run with
\texttt{python run\_all\_audits.py}.  The release environment used for the
final audit was Python~3.13.5, NetworkX~3.6.1, and pdf\TeX{} 1.40.26
(TeX Live 2025/dev).  A permanent repository or archival DOI, once assigned,
should point to the checksum-verified supplement distributed with this
manuscript.

\section*{Acknowledgements}

The author thanks the developers and maintainers of NetworkX for the graph
enumeration and verification tools used in the computational audit.

\begin{thebibliography}{99}

\bibitem{AIM2008}
AIM Minimum Rank -- Special Graphs Work Group,
\emph{Zero forcing sets and the minimum rank of graphs},
Linear Algebra Appl. \textbf{428} (2008), 1628--1648.

\bibitem{BoyerEtAl2019}
K. Boyer, B. Brimkov, S. English, D. Ferrero, A. Keller, R. Kirsch,
M. Phillips, and C. Reinhart,
\emph{The zero forcing polynomial of a graph},
Discrete Appl. Math. \textbf{258} (2019), 35--48.

\bibitem{CameronEtAl2026}
B. Cameron, J. Janssen, R. Mathew, and Z. Zhang,
\emph{An approximation algorithm for zero forcing},
Discrete Appl. Math. \textbf{383} (2026), 227--242.

\bibitem{CurtisEtAl2024}
B. Curtis, L. Gan, J. Haddock, R. Lawrence, and S. Spiro,
\emph{Zero forcing with random sets},
Discrete Math. \textbf{347} (2024), Article 113944.

\bibitem{AbiadGhasemiNezhad2026}
A. Abiad and S. Ghasemi Nezhad,
\emph{Fort abundance in zero forcing},
arXiv:2609.11754, 2026.

\bibitem{BeckerEtAl2025}
P. Becker, T. R. Cameron, D. Hanely, B. Ong, and J. P. Previte,
\emph{On the number of minimal forts of a graph},
Graphs Combin. \textbf{41} (2025), Article 25,
DOI: \url{https://doi.org/10.1007/s00373-025-02891-z}.

\bibitem{CameronLi2025}
T. R. Cameron and K. Li,
\emph{On the minimal forts of trees},
arXiv:2512.12874, 2025.

\bibitem{DatKenter2026}
N. H. Dat and F. H. J. Kenter,
\emph{On the number of zero forcing minimal forts on trees},
arXiv:2605.07298, 2026.

\bibitem{BaurStrassen1983}
W. Baur and V. Strassen,
\emph{The complexity of partial derivatives},
Theoret. Comput. Sci. \textbf{22} (1983), 317--330.

\bibitem{BhyravarapuEtAl2025}
S. Bhyravarapu, L. Kanesh, M. Kundu, D. Lokshtanov, and S. Saurabh,
\emph{Parameterized algorithms for Power Edge Set and Zero Forcing Set},
in \emph{Combinatorial Algorithms (IWOCA 2025)}, Lecture Notes in Computer
Science \textbf{15885}, 349--361, Springer, 2025,
DOI: \url{https://doi.org/10.1007/978-3-031-98740-3_25}.

\bibitem{GoettlicherEtAl2026}
M. G{\"o}ttlicher, L. Hofmann, and C. Niederbudde,
\emph{From FPT to W[P]: Classifying Zero Forcing, Power Domination and Their Variants},
arXiv:2609.29966, 2026.

\bibitem{McKayRayValliere2026}
A. McKay, A. Ray, and J. Valliere,
\emph{Maximizing the discrepancy between zero forcing parameters relative to graph order},
arXiv:2609.35522, 2026.

\bibitem{German2026}
S. German,
\emph{The path-extremal conjecture for zero forcing: distance-hereditary graphs and a split-decomposition reduction},
in \emph{Computing and Combinatorics (COCOON 2026)}, 409--422,
Springer, 2026, DOI: \url{https://doi.org/10.1007/978-981-92-3309-0_31}.

\bibitem{MenonSingh2025}
K. Menon and A. Singh,
\emph{Exploring the influence of graph operations on zero forcing sets},
Discrete Math. \textbf{348} (2025), Article 114516.

\bibitem{Jacob2025}
B. C. Jacob,
\emph{Well-failed graphs},
Australasian J. Combin. \textbf{93} (2025), 359--374.

\bibitem{ShamsizadehEtAl2021}
M. Shamsizadeh, M. M. Zahedi, M. Golmohamadian, and K. Abolpour,
\emph{Zero-forcing finite automata},
Int. J. Ind. Math. \textbf{13} (2021), 477--488.

\bibitem{ShamsizadehEtAl2024}
M. Shamsizadeh, M. M. Zahedi, K. Abolpour, and M. De la Sen,
\emph{Recognizable languages of $k$-forcing automata},
Math. Comput. Appl. \textbf{29} (2024), Article 32.

\bibitem{Wang2026}
X. Wang,
\emph{Modular compression of leaky forts and leaky forcing polynomials},
preprint, May 2026, DOI: 10.5281/zenodo.20388397.

\bibitem{AbiadMoghaddas2025}
A. Abiad and M. Moghaddas,
\emph{On the $Z_q$-forcing number: computational approach and exact values},
J. Phys. Complexity \textbf{6} (2025), Article 045005.

\bibitem{ArnborgLagergrenSeese1991}
S. Arnborg, J. Lagergren, and D. Seese,
\emph{Easy problems for tree-decomposable graphs},
J. Algorithms \textbf{12} (1991), 308--340.

\bibitem{BrimkovFastHicks2019}
B. Brimkov, C. C. Fast, and I. V. Hicks,
\emph{Computational approaches for zero forcing and related problems},
European J. Oper. Res. \textbf{273} (2019), 889--903.

\bibitem{CourcelleMakowskyRotics2001}
B. Courcelle, J. A. Makowsky, and U. Rotics,
\emph{On the fixed parameter complexity of graph enumeration problems definable in monadic second-order logic},
Discrete Appl. Math. \textbf{108} (2001), 23--52.

\bibitem{Gurski2017}
F. Gurski,
\emph{The behavior of clique-width under graph operations and graph transformations},
Theory Comput. Syst. \textbf{60} (2017), 346--376.

\bibitem{Lovasz2006}
L. Lov\'asz,
\emph{The rank of connection matrices and the dimension of graph algebras},
European J. Combin. \textbf{27} (2006), 962--970.

\bibitem{KotekMakowsky2014}
T. Kotek and J. A. Makowsky,
\emph{Connection matrices and the definability of graph parameters},
Logical Methods in Computer Science \textbf{10} (2014), no.~4, 1--33.

\bibitem{CameronHogbenEtAl2026}
T. R. Cameron, L. Hogben, F. H. J. Kenter, S. A. Mojallal, and H. Schuerger,
\emph{Forts, (fractional) zero forcing, and Cartesian products of graphs},
Australasian J. Combin. \textbf{95} (2026), 214--247.

\bibitem{Row2012}
D. D. Row,
\emph{A technique for computing the zero forcing number of a graph with a cut-vertex},
Linear Algebra Appl. \textbf{436} (2012), 4423--4432.

\end{thebibliography}

\end{document}
